\chapter{Roots of Unity Everywhere}

\begin{center}\itshape
God made the integers, all else is the work of man.\\
\upshape --- Leopold Kronecker, as reported by Heinrich Weber (1893)
\end{center}

Write $x^n-1$ as the product of its irreducible factors over $\Q$. For $n=1,2,\dots,104$ every coefficient of every factor
that appears is $-1$, $0$ or $1$. The factors are the cyclotomic polynomials $\Phi_d$, and after a hundred of them anyone
would bet that this is a theorem. In 1883 the Austrian mathematician Adalbert Migotti proved that it is a theorem when $n$
is a product of two primes, and for a while that seemed to be the whole story. Then comes $\Phi_{105}$, of degree $48$, the
first cyclotomic polynomial whose index has three odd prime factors. Its coefficient of $x^7$ is $-2$ (Figure~\ref{fig:phi105}).
Issai Schur showed in a letter to Edmund Landau that the coefficients are in fact unbounded, and in 1987 Jiro Suzuki proved
that \emph{every} integer is a coefficient of some $\Phi_n$.

So the roots of unity, the most symmetric numbers there are, carry some arithmetic surprises. This chapter collects them.
The roots of unity filter counts subsets, dice and tilings. Cyclotomic polynomials hide units, discriminants and quadratic
reciprocity. And the same Fourier modes that diagonalise a circulant matrix explain why polygons become regular when you
keep taking midpoints, and why a random walk on the fifth roots of unity, although it stays in the plane, escapes to
infinity.

\begin{think}
Before reading on, compute the coefficient of $x^7$ in $\Phi_{105}$ by hand. You only need the first eight coefficients, so
you may throw away every power $x^d$ with $d\ge8$. Which factors $1-x^d$ survive, and with which exponents?
\end{think}

\begin{figure}[ht]
\centering
\fig{ch02-phi105}
\caption{The $49$ coefficients of $\Phi_{105}$, from $x^0$ to $x^{48}$. Blue bars are $+1$, yellow bars $-1$. The two red bars, at $x^7$
and at $x^{41}$, are the first coefficients equal to $-2$ in the whole cyclotomic family (the polynomial is palindromic, so
they come in pairs).}
\label{fig:phi105}
\end{figure}

\section{Filters made of roots}

If $\omega=e^{2i\pi/n}$, then $\frac1n\sum_{j=0}^{n-1}\omega^{jk}$ is $1$ when $n\mid k$ and $0$ otherwise. This one line
turns ``is divisible by $n$'' into an average, and averages are easy to compute when the thing averaged is a product. The
first problems all run on this engine.

\begin{exercise}[$\star\star$]\label{pb:subsetsum}
\begin{enumerate}[label=(\alph*)]
\item Let $p$ be an odd prime. How many subsets of $\{1,2,\dots,p\}$ (the empty set included) have a sum divisible by $p$?
\item How many $p$-element subsets of $\{1,2,\dots,2p\}$ have a sum divisible by $p$?
\item Let $n\ge1$ and $0\le k\le n$, and let $N(n,k)$ be the number of $k$-element subsets of $\Z/n\Z$ whose sum is $0$.
Show that
\[
N(n,k)=\frac1n\sum_{d\mid\gcd(n,k)}(-1)^{k+k/d}\,\varphi(d)\binom{n/d}{k/d}.
\]
\end{enumerate}
\end{exercise}

\begin{exercise}[$\star$]\label{pb:dicemod}
Throw $n\ge1$ fair dice and let $S_n$ be the sum.
\begin{enumerate}[label=(\alph*)]
\item Show that $P(7\mid S_n)=\frac17\bigl(1+6(-\frac16)^n\bigr)$.
\item Show that $P(6\mid S_n)=\frac16$ exactly, first with roots of unity, then with no computation at all.
\item Compute $P(5\mid S_n)$. For which $n$ is it larger than $\frac15$?
\end{enumerate}
\end{exercise}

\begin{exercise}[$\star\star$]\label{pb:ramsum}
For integers $q\ge1$ and $n$, the \emph{Ramanujan sum} is $c_q(n)=\sum e^{2i\pi an/q}$, over $1\le a\le q$ with
$\gcd(a,q)=1$.
\begin{enumerate}[label=(\alph*)]
\item Show that $c_q(n)=\sum_{d\mid\gcd(q,n)}\mu(q/d)\,d$. In particular $c_q(n)$ is an integer.
\item Show Hölder's formula: with $g=\gcd(q,n)$,
\[
c_q(n)=\mu\Bigl(\frac qg\Bigr)\frac{\varphi(q)}{\varphi(q/g)}.
\]
\item Let $q$ and $r$ divide $N$. Show that $\frac1N\sum_{n=1}^Nc_q(n)c_r(n)$ is $\varphi(q)$ if $q=r$ and $0$ otherwise.
\end{enumerate}
\end{exercise}

\begin{exercise}[$\star\star$]\label{pb:bricks}
\begin{enumerate}[label=(\alph*)]
\item An $a\times b$ rectangle of unit squares is tiled by $1\times k$ bricks, each placed horizontally or vertically.
Show that $k$ divides $a$ or $b$. Can a $10\times10$ board be tiled by $1\times4$ bricks?
\item A rectangle of the plane is tiled by finitely many smaller rectangles with sides parallel to its own. Each small
rectangle has at least one side of integer length. Show that the big rectangle has a side of integer length.
\end{enumerate}
\end{exercise}

\begin{exercise}[$\star\star$]\label{pb:prouhet}
Let $s(n)$ be the number of ones in the binary expansion of $n$, and $t_n=(-1)^{s(n)}$.
\begin{enumerate}[label=(\alph*)]
\item Let $N\ge1$. Show that $\sum_{n=0}^{2^N-1}t_n\,n^k=0$ for $k=0,1,\dots,N-1$.
\item Show that for $k=N$ the sum equals $(-1)^NN!\,2^{N(N-1)/2}$.
\item Split $\{0,1,\dots,15\}$ into two sets of eight numbers with the same sum, the same sum of squares and the same sum
of cubes.
\end{enumerate}
\end{exercise}

\begin{exercise}[$\star\star$]\label{pb:lightsout}
Lamps numbered by $\Z/n\Z$, $n\ge3$, stand on a circle, each on or off. Pressing lamp $i$ switches the state of the lamps
$i-1$, $i$ and $i+1$.
\begin{enumerate}[label=(\alph*)]
\item Show that every configuration can be switched off if and only if $3\nmid n$.
\item Let $3\mid n$. For $r\in\{0,1,2\}$, let $S_r$ be the number of lit lamps $i$ with $i\equiv r\pmod3$. Show that a
configuration can be switched off if and only if $S_0$, $S_1$ and $S_2$ have the same parity. What proportion of the
configurations is that?
\end{enumerate}
\end{exercise}

\begin{exercise}[$\star\star\star$]\label{pb:sieving}
\emph{Cyclic sieving.} Let $[m]_q=1+q+\dots+q^{m-1}$, $[m]_q!=[1]_q[2]_q\cdots[m]_q$, and for $0\le k\le n$
\[
\qbinom nk=\frac{[n]_q!}{[k]_q!\,[n-k]_q!}.
\]
\begin{enumerate}[label=(\alph*)]
\item Show that $\qbinom nk=\qbinom{n-1}{k-1}+q^k\qbinom{n-1}k$, and deduce that $\qbinom nk$ is a polynomial in $q$ with
non-negative integer coefficients.
\item Let $m\mid n$ and let $\zeta$ be a primitive $m$-th root of unity. Show that the value of this polynomial at $\zeta$
is $\binom{n/m}{k/m}$ if $m\mid k$, and $0$ otherwise.
\item The rotation $g:x\mapsto x+1$ of $\Z/n\Z$ acts on the $k$-element subsets of $\Z/n\Z$. Show that for every $j$, the
number of subsets fixed by $g^j$ is the value of $\qbinom nk$ at $q=e^{2i\pi j/n}$.
\item Deduce the number of necklaces with $n$ beads, $k$ of them black, up to rotation, and check that there are $4$ for
$n=6$, $k=3$.
\end{enumerate}
\end{exercise}

\section{Cyclotomic polynomials up close}

Recall that $X^n-1=\prod_{d\mid n}\Phi_d$, that $\Phi_n\in\Z[X]$ is monic of degree $\varphi(n)$, and that its roots are the
primitive $n$-th roots of unity (Tome~1, Chapter~5). The next problems look at these polynomials from every side: their
values, their coefficients, their factors modulo $p$, and the numbers they make.

\begin{exercise}[$\star\star$]\label{pb:phipm1}
\begin{enumerate}[label=(\alph*)]
\item Show that $\Phi_n(1)=e^{\Lambda(n)}$ for $n\ge2$: it is $p$ if $n$ is a power of a prime $p$, and $1$ otherwise.
\item Compute $\Phi_n(-1)$ for every $n\ge3$.
\item Let $n\ge2$ not be a prime power and $\zeta=e^{2i\pi/n}$. Show that $1-\zeta$ is invertible in the ring $\Z[\zeta]$.
\end{enumerate}
\end{exercise}

\begin{exercise}[$\star\star\star$]\label{pb:phipq}
\begin{enumerate}[label=(\alph*)]
\item Let $p<q$ be primes. Show that every coefficient of $\Phi_{pq}$ is $-1$, $0$ or $1$, and that the non-zero
coefficients alternate between $1$ and $-1$.
\item Show that the coefficient of $x^7$ in $\Phi_{105}$ is $-2$.
\end{enumerate}
\end{exercise}

\begin{exercise}[$\star\star$]\label{pb:vanish}
\begin{enumerate}[label=(\alph*)]
\item Let $p$ be prime and $\zeta=e^{2i\pi/p}$. Show that rationals $a_0,\dots,a_{p-1}$ satisfy $\sum_ka_k\zeta^k=0$ if and
only if they are all equal. Which sets of $p$-th roots of unity have sum $0$?
\item Find all non-empty sets of sixth roots of unity with sum $0$. (There are $9$.)
\item Let $n=p^a$ be a prime power. Show that a set of $n$-th roots of unity has sum $0$ if and only if it is a union of
sets of the form $\{\eta,\eta\omega,\dots,\eta\omega^{p-1}\}$, with $\omega=e^{2i\pi/p}$.
\end{enumerate}
\end{exercise}

\begin{exercise}[$\star\star\star$]\label{pb:kronecker}
\emph{Kronecker's theorem.}
\begin{enumerate}[label=(\alph*)]
\item Let $P\in\Z[X]$ be monic with $P(0)\ne0$, and suppose that all its complex roots have modulus at most $1$. Show
that they are all roots of unity, and that if $P$ is irreducible, $P$ is a cyclotomic polynomial.
\item Let $P\in\Z[X]$ be monic, with all its roots real and in $[-2,2]$. Show that its roots are of the form $2\cos(2\pi r)$
with $r\in\Q$.
\item The number $z=\frac{3+4i}5$ has modulus $1$, and so does its only conjugate $\bar z$. Show that it is not a root of
unity. Which hypothesis of (a) fails?
\end{enumerate}
\end{exercise}

\begin{exercise}[$\star\star$]\label{pb:primsum}
Let $p$ be an odd prime.
\begin{enumerate}[label=(\alph*)]
\item Show that the sum of the primitive roots modulo $p$ in $\{1,\dots,p-1\}$ is congruent to $\mu(p-1)$ modulo $p$.
\item Show that their product is congruent to $1$ modulo $p$ when $p>3$.
\item More generally, show that the sum of the $k$-th powers of the primitive roots is congruent modulo $p$ to the
Ramanujan sum $c_{p-1}(k)$ of Problem~\ref{pb:ramsum}.
\end{enumerate}
\end{exercise}

\begin{exercise}[$\star\star$]\label{pb:phimodp}
\begin{enumerate}[label=(\alph*)]
\item Let $p$ be a prime not dividing $n$, and $f=\operatorname{ord}_n(p)$. Show that $\Phi_n$ modulo $p$ is a product of
$\varphi(n)/f$ distinct irreducible polynomials of $\mathbb F_p[X]$, all of degree $f$.
\item Deduce that $X^4+1$ is irreducible over $\Q$ but factors modulo every prime.
\item Let $n=p^am$ with $a\ge1$ and $p\nmid m$. Show that $\Phi_n\equiv\Phi_m^{\varphi(p^a)}\pmod p$.
\end{enumerate}
\end{exercise}

\section{Rationality and radicals}

Roots of unity are algebraic, so they can be rational only in the simplest cases, and they can be written with radicals.
The problems of this section turn both facts into statements about angles, cube roots, and quadratic reciprocity.

\begin{exercise}[$\star\star$]\label{pb:rattan}
\begin{enumerate}[label=(\alph*)]
\item Let $\zeta_n=e^{2i\pi/n}$. Show that the roots of unity in $\Q(\zeta_n)$ are the $n$-th roots of unity if $n$ is
even, and the $2n$-th roots of unity if $n$ is odd.
\item Let $\theta$ be a rational multiple of $\pi$ with $\tan\theta\in\Q$. Show that $\tan\theta\in\{-1,0,1\}$.
\item Show that the only regular polygon whose vertices all lie in $\Z^2$ is the square.
\item (Niven) Let $\theta$ be a rational multiple of $\pi$ with $\cos\theta\in\Q$. Show that
$\cos\theta\in\{0,\pm\frac12,\pm1\}$.
\end{enumerate}
\end{exercise}

\begin{exercise}[$\star\star\star$]\label{pb:ramacube}
\emph{Ramanujan's cube roots.} With real cube roots, show that
\[
\sqrt[3]{\cos\tfrac{2\pi}7}+\sqrt[3]{\cos\tfrac{4\pi}7}+\sqrt[3]{\cos\tfrac{8\pi}7}=\sqrt[3]{\frac{5-3\sqrt[3]7}2}.
\]
\end{exercise}

\begin{exercise}[$\star\star\star$]\label{pb:discphi}
Let $p$ be an odd prime, $\zeta=e^{2i\pi/p}$ and $p^*=(-1)^{(p-1)/2}p$.
\begin{enumerate}[label=(\alph*)]
\item Show that the discriminant $\prod_{1\le i<j\le p-1}(\zeta^i-\zeta^j)^2$ of $\Phi_p$ equals $(p^*)\,p^{p-3}$, and
deduce that $\sqrt{p^*}\in\Q(\zeta)$.
\item (Zolotarev) For $a$ prime to $p$, let $\pi_a$ be the permutation $x\mapsto ax$ of $\Z/p\Z$. Show that the sign of
$\pi_a$ is the Legendre symbol $\bigl(\frac ap\bigr)$.
\item Let $\sigma_a$ be the automorphism of $\Q(\zeta)$ with $\sigma_a(\zeta)=\zeta^a$. Show that
$\sigma_a(\sqrt{p^*})=\bigl(\frac ap\bigr)\sqrt{p^*}$.
\end{enumerate}
\end{exercise}

\begin{exercise}[Journey]\label{pb:eisensine}
\emph{Eisenstein's trigonometric proof of quadratic reciprocity.}
\begin{enumerate}[label=(\alph*)]
\item Let $q$ be odd. Show that for every real $x$ that is not a multiple of $\pi$,
\[
\frac{\sin qx}{\sin x}=(-4)^{(q-1)/2}\prod_{k=1}^{(q-1)/2}\Bigl(\sin^2x-\sin^2\frac{2\pi k}q\Bigr).
\]
\item Let $p$ be an odd prime and $q$ an integer prime to $p$. Show that
\[
\Bigl(\frac qp\Bigr)=\prod_{j=1}^{(p-1)/2}\frac{\sin(2\pi qj/p)}{\sin(2\pi j/p)}.
\]
\item Deduce the law of quadratic reciprocity: for distinct odd primes $p$ and $q$,
$\bigl(\frac qp\bigr)\bigl(\frac pq\bigr)=(-1)^{\frac{p-1}2\frac{q-1}2}$.
\end{enumerate}
\end{exercise}

\section{Polygons and walks}

Number the vertices of a polygon by $\Z/n\Z$. Any operation that treats all vertices alike (take midpoints, erect triangles
on the sides) is a circulant linear map, and the discrete Fourier transform of Tome~4, Chapter~7 diagonalises it: the
eigenvectors are the regular polygons $k\mapsto\omega^{jk}$. Geometry then reduces to reading off eigenvalues.

\begin{exercise}[$\star\star$]\label{pb:jensen}
\begin{enumerate}[label=(\alph*)]
\item Show that for real $x$ and $\theta$ and $n\ge1$,
\[
\prod_{k=0}^{n-1}\Bigl(x^2-2x\cos\Bigl(\theta+\frac{2\pi k}n\Bigr)+1\Bigr)=x^{2n}-2x^n\cos n\theta+1.
\]
\item Deduce Poisson's integral: for $|x|\ne1$, $\int_0^\pi\log(1-2x\cos\theta+x^2)\dd\theta$ is $0$ if $|x|<1$ and
$2\pi\log|x|$ if $|x|>1$.
\item Let $P=a\prod_{k=1}^d(X-z_k)\in\C[X]$, $a\ne0$, with no root of modulus $1$. Show that
\[
\frac1{2\pi}\int_0^{2\pi}\log\bigl|P(e^{i\theta})\bigr|\dd\theta=\log|a|+\sum_{k=1}^d\log\max(1,|z_k|).
\]
What does Problem~\ref{pb:kronecker} say about the polynomials of $\Z[X]$ for which this is $0$?
\end{enumerate}
\end{exercise}

\begin{exercise}[$\star\star$]\label{pb:midpoint}
Let $n\ge3$ and $z=(z_0,\dots,z_{n-1})\in\C^n$ be the vertices of a polygon with centroid $0$. Its \emph{midpoint
polygon} is $Mz$, with $(Mz)_k=\frac12(z_k+z_{k+1})$, indices modulo $n$.
\begin{enumerate}[label=(\alph*)]
\item Show that there are complex numbers $a$ and $b$, depending on $z$, such that, with $\omega=e^{2i\pi/n}$,
\[
\frac{(M^tz)_k}{\cos^t(\pi/n)}-\bigl(a\,e^{i\pi t/n}\omega^k+b\,e^{-i\pi t/n}\omega^{-k}\bigr)\xrightarrow[t\to\infty]{}0
\quad\text{for each }k.
\]
\item Show that $k\mapsto a\omega^k+b\omega^{-k}$ is the image of the regular polygon $k\mapsto\omega^k$ by a real-linear
map of the plane, invertible when $|a|\ne|b|$. So the iterated midpoint polygons, rescaled, become \emph{affinely regular}
(Figure~\ref{fig:midpoint}).
\item What happens for $n=4$ after a single step?
\end{enumerate}
\end{exercise}

\begin{figure}[ht]
\centering
\fig{ch02-midpoint}
\caption{A random heptagon and its iterated midpoint polygons, each rescaled to the same size. The modes of
Problem~\ref{pb:midpoint} die at different speeds, and what survives is an affine image of the regular heptagon: an ellipse
with seven equally spaced points (in the sense of the parametrisation) on it.}
\label{fig:midpoint}
\end{figure}

\begin{exercise}[$\star\star$]\label{pb:napoleon}
\emph{Napoleon's theorem.} Let $ABC$ be a triangle, labelled counterclockwise. Erect equilateral triangles outwards on its
three sides, and let $P$, $Q$, $R$ be their centres.
\begin{enumerate}[label=(\alph*)]
\item Show that $PQR$ is equilateral, with the same centroid as $ABC$.
\item Do the same with equilateral triangles erected inwards, with centres $P'$, $Q'$, $R'$. Show that $P'Q'R'$ is
equilateral, and that the area of $PQR$ minus the area of $P'Q'R'$ equals the area of $ABC$.
\end{enumerate}
\end{exercise}

\begin{exercise}[$\star\star\star$]\label{pb:pdn}
\emph{The Petr--Douglas--Neumann theorem.} Let $n\ge3$, $z\in\C^n$ a polygon, and $\theta\in\left]0,2\pi\right[$. Let
$A_\theta z$ be the polygon whose $k$-th vertex is the point $p$ with $z_k-p=e^{i\theta}(z_{k+1}-p)$: the apex of the
isosceles triangle with apex angle $\theta$ erected on the side $z_kz_{k+1}$.
\begin{enumerate}[label=(\alph*)]
\item Find the eigenvalues of the linear map $A_\theta$ on the Fourier modes $k\mapsto\omega^{jk}$.
\item Show that applying $A_{2\pi k/n}$ for $k=1,2,\dots,n-2$, in any order, gives a regular $n$-gon (possibly reduced to a
point) whose centre is the centroid of $z$ (Figure~\ref{fig:pdn}).
\item What do $n=3$ and $n=4$ say?
\end{enumerate}
\end{exercise}

\begin{figure}[ht]
\centering
\fig{ch02-pdn}
\caption{The Petr--Douglas--Neumann theorem for a pentagon. On each side erect isosceles triangles with apex angle $72^\circ$
(dashed), keep the apexes, repeat with $144^\circ$ and then $216^\circ$ (the last triangles point inwards). Each step kills
one Fourier mode, and after three steps only a regular pentagon is left.}
\label{fig:pdn}
\end{figure}

\begin{exercise}[$\star\star\star$]\label{pb:rootwalk}
Let $n\ge3$. A walker starts at $0\in\C$ and makes independent steps $\xi_1,\xi_2,\dots$, each uniform on the $n$-th roots
of unity. Let $S_m=\xi_1+\dots+\xi_m$.
\begin{enumerate}[label=(\alph*)]
\item Show that $E|S_m|^2=m$ and $E|S_m|^4=2m^2-m$, whatever $n\ge3$.
\item Let $n=p$ be prime. Show that $S_m=0$ only if $p\mid m$, and that
\[
P(S_{pN}=0)=\frac{(pN)!}{N!^p\,p^{pN}}\sim\frac{\sqrt p}{(2\pi N)^{(p-1)/2}}.
\]
\item Deduce that for $p=3$ the walker returns to $0$ infinitely often almost surely, while for $p\ge5$, although it never
leaves the plane, it returns only finitely many times.
\end{enumerate}
\end{exercise}

\section{Hints}

\begin{hint}{pb:subsetsum}
Expand $\prod_{a\in\Z/n\Z}(1+y\,\omega^{ja})$ and filter with $\frac1n\sum_j$. For a root of unity $\eta$ of order $d$,
$\prod_{\varepsilon^d=1}(1+y\varepsilon)=1-(-y)^d$.
\end{hint}

\begin{hint}{pb:dicemod}
$P(m\mid S_n)=\frac1m\sum_{j}E[\omega^{jS_n}]$ and $E[\omega^{jS_n}]=\bigl(\frac16(\omega^j+\dots+\omega^{6j})\bigr)^n$. For
(b) without computation, condition on the first $n-1$ dice.
\end{hint}

\begin{hint}{pb:ramsum}
(a) Sort the $a\in\{1,\dots,q\}$ by $\gcd(a,q)$ and use Möbius inversion. (b) Both sides are multiplicative in $q$; check
prime powers. (c) Expand both sums and use that $a/q$ and $b/r$, in lowest terms, are equal modulo $1$ only if they are
the same fraction.
\end{hint}

\begin{hint}{pb:bricks}
(a) Put $\zeta^{i+j}$, $\zeta=e^{2i\pi/k}$, in the cell $(i,j)$. (b) Integrate $e^{2i\pi(x+y)}$ over each rectangle.
\end{hint}

\begin{hint}{pb:prouhet}
$\sum_{n<2^N}t_nx^n=\prod_{j<N}(1-x^{2^j})$ has a zero of order $N$ at $x=1$. Apply $(x\frac d{dx})^k$ and evaluate at $1$.
\end{hint}

\begin{hint}{pb:lightsout}
Identify configurations with $\mathbb F_2[x]/(x^n-1)$: pressing the set $u$ produces $(x^{-1}+1+x)u$. For (b), the matrix is
symmetric, so its image is the orthogonal of its kernel.
\end{hint}

\begin{hint}{pb:sieving}
(b) Use $\qbinom nk=\prod_{i=1}^k\frac{1-q^{n-k+i}}{1-q^i}$ and compute the limit as $q\to\zeta$: count the factors that
vanish at $\zeta$ above and below. (c) A subset fixed by $g^j$ is a union of orbits of $\langle g^j\rangle$. (d) Burnside.
\end{hint}

\begin{hint}{pb:phipm1}
(a) Evaluate $\frac{x^n-1}{x-1}=\prod_{d\mid n,d>1}\Phi_d$ at $1$ and take logarithms. (b) Prove $\Phi_{2m}(x)=\Phi_m(-x)$
for odd $m>1$, and $\Phi_{pm}(x)=\Phi_m(x^p)$ when $p\mid m$. (c) $\Phi_n(1)$ is a product of factors $1-\zeta^a$.
\end{hint}

\begin{hint}{pb:phipq}
(a) $\Phi_{pq}(x)=\frac{(1-x)(1-x^{pq})}{(1-x^p)(1-x^q)}$, and $\deg\Phi_{pq}<pq$: compute modulo $x^{pq}$. How many ways
can $m<pq$ be written as $ip+jq$ with $i,j\ge0$? (b) Use $\Phi_n=\prod_{d\mid n}(1-x^d)^{\mu(n/d)}$ modulo $x^8$.
\end{hint}

\begin{hint}{pb:vanish}
(a) $\Phi_p$ is irreducible. (b) Write everything in the basis $1,\zeta$ of $\Q(\zeta_6)$, using $\zeta^3=-1$. (c) Group the
exponents by their residue modulo $p^{a-1}$; over $\Q(\omega)$, $\zeta_{p^a}$ has degree $p^{a-1}$.
\end{hint}

\begin{hint}{pb:kronecker}
(a) If $C$ is the companion matrix of $P$, the characteristic polynomial of $C^k$ has integer coefficients and roots
$\alpha_i^k$. There are only finitely many monic integer polynomials of degree $d$ with all roots in the closed unit disc.
(b) Look at $x^dP(x+\frac1x)$.
\end{hint}

\begin{hint}{pb:primsum}
In $\mathbb F_p$, the primitive roots are the $\gamma^a$ with $\gcd(a,p-1)=1$, for one primitive root $\gamma$. Copy the
proof of Problem~\ref{pb:ramsum}(a) with $\gamma$ in place of $e^{2i\pi/q}$.
\end{hint}

\begin{hint}{pb:phimodp}
(a) $X^n-1$ has distinct roots modulo $p$, so the roots of $\Phi_n$ modulo $p$ have order exactly $n$. An element of order $n$
lies in $\mathbb F_{p^f}$ exactly when $n\mid p^f-1$. (c) Use $\Phi_m(x^p)\equiv\Phi_m(x)^p\pmod p$.
\end{hint}

\begin{hint}{pb:rattan}
(a) If $\zeta_m\in\Q(\zeta_n)$, then $\zeta_{\operatorname{lcm}(m,n)}\in\Q(\zeta_n)$; compare $\varphi$. (b)
$\frac{1+it}{1-it}=e^{2i\theta}$. (d) $2\cos\theta$ is an algebraic integer.
\end{hint}

\begin{hint}{pb:ramacube}
Let $a,b,c$ be the three cube roots, $s=a+b+c$, $t=ab+bc+ca$, and use that the cosines are the roots of
$8x^3+4x^2-4x-1$. Express $s^3$ and $t^3$ in terms of $m=st$, and find a cubic equation for $m$.
\end{hint}

\begin{hint}{pb:discphi}
(a) $\prod_{i\ne j}(\alpha_i-\alpha_j)=\prod_iP'(\alpha_i)$, and $(x-1)\Phi_p(x)=x^p-1$. (b) $a\mapsto\operatorname{sgn}\pi_a$
is a homomorphism; what is the cycle type of $\pi_g$ for a primitive root $g$? (c) Apply $\sigma_a$ to
$\prod_{i<j}(\zeta^j-\zeta^i)$.
\end{hint}

\begin{hint}{pb:eisensine}
(a) $\frac{\sin qx}{\sin x}$ is a polynomial in $\sin^2x$ of degree $\frac{q-1}2$; find its leading coefficient and its
roots. (b) Gauss's lemma: write $qj\equiv\pm r_j$ with $1\le r_j\le\frac{p-1}2$. (c) Plug (a) into (b) and exchange $p$
and $q$.
\end{hint}

\begin{hint}{pb:jensen}
(a) $x^2-2x\cos\varphi+1=(x-e^{i\varphi})(x-e^{-i\varphi})$. (b) Riemann sums at the points $\frac{2\pi k}n$. (c) Apply (b)
to each factor.
\end{hint}

\begin{hint}{pb:midpoint}
Expand $z$ in the discrete Fourier basis: $M$ multiplies the mode $\omega^{jk}$ by $\frac{1+\omega^j}2$.
\end{hint}

\begin{hint}{pb:napoleon}
The centre of an equilateral triangle on a side is the apex of an isosceles triangle with apex angle $120^\circ$. Use the
modes of $\Z/3\Z$ and the area formula $\frac12\operatorname{Im}\sum_k\bar z_kz_{k+1}$.
\end{hint}

\begin{hint}{pb:pdn}
Solve for $p$, and apply $A_\theta$ to $z_k=\omega^{jk}$. Which mode does $A_{2\pi k/n}$ kill?
\end{hint}

\begin{hint}{pb:rootwalk}
(a) $E[\xi^r\bar\xi^s]=0$ unless $n\mid r-s$. (b) $\sum_ka_k\zeta^k=0$ with $a_k\ge0$ integers forces all $a_k$ equal
(Problem~\ref{pb:vanish}). Stirling. (c) The recurrence criterion of Tome~3, Chapter~7.
\end{hint}

\section{Solutions}

\begin{solution}{pb:subsetsum}
\idea ``The sum is divisible by $n$'' is a filter: average $\omega^{j\cdot\text{sum}}$ over $j$. The generating function
of subsets is a product, $\prod_a(1+y\,x^a)$, and at $x=\omega^j$ the product runs over the powers of a root of unity of
order $d$, each repeated $\frac nd$ times. Such products collapse to $1-(-y)^d$.

\writeup Fix $n$ and let $\omega=e^{2i\pi/n}$. For a subset $A$ of $\Z/n\Z$, $\frac1n\sum_{j=0}^{n-1}\omega^{j\sum A}$ is $1$
if $\sum A=0$ and $0$ otherwise (Tome~0, Chapter~2). So
\[
N(n,k)=\frac1n\sum_{j=0}^{n-1}[y^k]F_j(y),\qquad F_j(y)=\prod_{a\in\Z/n\Z}(1+y\,\omega^{ja}),
\]
because expanding $F_j$ gives $\sum_Ay^{|A|}\omega^{j\sum A}$. Let $d$ be the order of $\omega^j$, that is,
$d=n/\gcd(n,j)$. As $a$ runs over $\Z/n\Z$, $\omega^{ja}$ runs over the $d$-th roots of unity, each $\frac nd$ times. Now
$\prod_{\varepsilon^d=1}(X-\varepsilon)=X^d-1$; put $X=-\frac1y$ and multiply by $(-y)^d$:
$\prod_\varepsilon(1+y\varepsilon)=1-(-y)^d$. Hence
\[
F_j(y)=\bigl(1-(-y)^d\bigr)^{n/d}=\sum_{i}\binom{n/d}i(-1)^i(-1)^{di}y^{di}.
\]
The coefficient of $y^k$ is $0$ unless $d\mid k$, and then it is $(-1)^{k/d+k}\binom{n/d}{k/d}$. Exactly $\varphi(d)$ values
of $j$ give the order $d$, for each $d\mid n$. Summing gives the formula of (c).

For (a), $\{1,\dots,p\}$ is $\Z/p\Z$, and we count subsets of all sizes, so we take $y=1$: $F_0(1)=2^p$ and, for $j\ne0$,
$d=p$ and $F_j(1)=1-(-1)^p=2$. The number is $\frac1p\bigl(2^p+2(p-1)\bigr)=\frac{2^p-2}p+2$. For $p=3$ this is $4$:
$\varnothing$, $\{3\}$, $\{1,2\}$, $\{1,2,3\}$.

For (b), $\{1,\dots,2p\}$ contains each residue twice, so the count is $\frac1p\sum_j[y^p]F_j(y)^2$. For $j=0$ this is
$[y^p](1+y)^{2p}=\binom{2p}p$; for $j\ne0$ it is $[y^p](1+y^p)^2=2$. The answer is $\frac1p\bigl(\binom{2p}p+2(p-1)\bigr)
=\frac{\binom{2p}p-2}p+2$. For $p=5$ this is $52$.
\tools Tome~0, Chapter~2 (roots of unity); Tome~3, Chapter~3 (generating functions).
\leads Part (b) is Problem~6 of the 1995 International Mathematical Olympiad, where most solutions were long counting
arguments. The formula of (c) also counts binary necklaces with a given number of black beads and a ``sum'' condition; the
same filter computes the number of solutions of $a_1+\dots+a_k=0$ in any finite abelian group, and it is the starting point
of the circle method of Hardy and Littlewood (Tome~13).
\end{solution}

\begin{solution}{pb:dicemod}
\idea Filter with $m$-th roots of unity. Everything rests on the value of $f(z)=z+z^2+\dots+z^6$ at the $m$-th roots of
unity, and for $m=5,6,7$ these sums of six consecutive powers are very simple.

\writeup Let $\omega=e^{2i\pi/m}$. As in Problem~\ref{pb:subsetsum},
$P(m\mid S_n)=\frac1m\sum_{j=0}^{m-1}E\bigl[\omega^{jS_n}\bigr]$, and by independence
$E[\omega^{jS_n}]=\bigl(f(\omega^j)/6\bigr)^n$.

(a) $m=7$. For $\zeta\ne1$ a $7$-th root of unity, $1+\zeta+\dots+\zeta^6=0$, so $f(\zeta)=-1$. Hence
$P(7\mid S_n)=\frac17\bigl(1+6(-\frac16)^n\bigr)$.

(b) $m=6$. For $\zeta\ne1$ a $6$-th root of unity, $f(\zeta)=\zeta(1+\dots+\zeta^5)=0$, so $P(6\mid S_n)=\frac16$. Without
computation: whatever the first $n-1$ dice show, exactly one of the six faces of the last die makes the total divisible by
$6$, so the conditional probability is $\frac16$, and so is the probability.

(c) $m=5$. For $\zeta\ne1$ a fifth root of unity, $f(\zeta)=(\zeta+\dots+\zeta^5)+\zeta^6=0+\zeta$. So
\[
P(5\mid S_n)=\frac15\Bigl(1+6^{-n}\sum_{j=1}^4\omega^{jn}\Bigr)=\frac15+\frac{5\cdot\mathbf 1_{5\mid n}-1}{5\cdot6^n},
\]
since $\sum_{j=1}^4\omega^{jn}$ is $4$ if $5\mid n$ and $-1$ otherwise. It exceeds $\frac15$ exactly when $5\mid n$. For
$n=1$ it is $\frac16$, for $n=2$ it is $\frac7{36}$, both below $\frac15$, as a direct count confirms.
\tools Tome~3, Chapters~2 and~3 (independence, generating functions); Tome~0, Chapter~2 (roots of unity).
\leads The error terms decay like $6^{-n}$ because $|f(\zeta)|$ is much smaller than $6$ for every $\zeta\ne1$. In general
$P(m\mid S_n)\to\frac1m$ for every $m$, and the speed is governed by $\max_{\zeta\ne1}|f(\zeta)|/6$: a spectral gap, the
same quantity that controls the mixing of random walks on groups (Tome~3, Chapter~8).
\end{solution}

\begin{solution}{pb:ramsum}
\idea The sum over \emph{all} $a$ is a geometric sum, $q$ or $0$. Sorting the $a$ by $\gcd(a,q)$ writes this full sum as a
sum of Ramanujan sums over the divisors of $q$, and Möbius inversion undoes it. For (b), multiplicativity reduces
everything to prime powers. For (c), think of $c_q(n)$ as a sum over the fractions with denominator exactly $q$.

\writeup Write $e(x)=e^{2i\pi x}$ and $\eta_q(n)=\sum_{a=1}^qe(an/q)$, which is $q$ if $q\mid n$ and $0$ otherwise.

(a) For $a\in\{1,\dots,q\}$ with $\gcd(a,q)=q/e$, write $a=\frac qeb$ with $1\le b\le e$ and $\gcd(b,e)=1$; then
$e(an/q)=e(bn/e)$. This is a bijection between $\{1,\dots,q\}$ and the pairs $(e,b)$ with $e\mid q$, $b$ prime to $e$. So
$\eta_q(n)=\sum_{e\mid q}c_e(n)$, and Möbius inversion (Tome~1, Chapter~4) gives
\[
c_q(n)=\sum_{e\mid q}\mu(q/e)\eta_e(n)=\sum_{e\mid q,\ e\mid n}\mu(q/e)\,e.
\]

(b) For fixed $n$, the function $h(e)=e\cdot\mathbf 1_{e\mid n}$ is multiplicative in $e$, so $c_q(n)=(\mu*h)(q)$ is
multiplicative in $q$ (Tome~1, Chapter~4). The right-hand side of Hölder's formula is multiplicative in $q$ too: if $q=q_1q_2$
with $\gcd(q_1,q_2)=1$, then $g=g_1g_2$ with $g_i=\gcd(q_i,n)$, and $\mu$ and $\varphi$ are multiplicative. So take $q=p^a$,
$a\ge1$, and $g=p^b$ with $b=\min(a,v_p(n))$. In (a) only $e=p^a$ and $e=p^{a-1}$ have $\mu(q/e)\ne0$:
\begin{itemize}
\item if $b=a$, $c_{p^a}(n)=p^a-p^{a-1}=\varphi(p^a)$;
\item if $b=a-1$, only $e=p^{a-1}$ divides $n$, and $c_{p^a}(n)=-p^{a-1}$;
\item if $b\le a-2$, neither divides $n$, and $c_{p^a}(n)=0$.
\end{itemize}
Hölder's formula gives $\mu(1)\varphi(p^a)$, then $\mu(p)\frac{\varphi(p^a)}{\varphi(p)}=-p^{a-1}$, then $0$, since
$\mu(p^{a-b})=0$ for $a-b\ge2$. They agree.

(c) The numbers $c_r(n)$ are integers, hence real, so $c_r(n)=\overline{c_r(n)}=\sum_be(-bn/r)$. Then
\[
\frac1N\sum_{n=1}^Nc_q(n)c_r(n)=\sum_{a,b}\frac1N\sum_{n=1}^Ne\Bigl(n\Bigl(\frac aq-\frac br\Bigr)\Bigr),
\]
with $a$ prime to $q$ and $b$ prime to $r$. Since $q,r\mid N$, $u=e(\frac aq-\frac br)$ is an $N$-th root of unity, and the
inner average is $1$ if $u=1$ and $0$ otherwise. And $u=1$ means $\frac aq\equiv\frac br\pmod1$; both fractions are in
lowest terms in $\left]0,1\right]$, so this means $q=r$ and $a=b$. The total is $\varphi(q)$ if $q=r$, and $0$ otherwise.
\tools Tome~1, Chapter~4 (Möbius inversion, multiplicative functions); Tome~0, Chapter~2 (roots of unity).
\leads Ramanujan (1918) expanded arithmetic functions in series of his sums, like Fourier series:
$\frac{\sigma(n)}n=\frac{\pi^2}6\sum_{q\ge1}\frac{c_q(n)}{q^2}$. Part (c) is the orthogonality
that makes such expansions work. Ramanujan sums also give the main terms in the circle method and in sums of the form
$\sum_n\Lambda(n)\Lambda(n+h)$ (Tome~13).
\end{solution}

\begin{solution}{pb:bricks}
\idea A colouring argument where the colours are roots of unity. Each brick covers $k$ consecutive powers of
$\zeta=e^{2i\pi/k}$, whose sum is $0$, so the whole board has weight $0$. But the board's weight factorises. The same idea
works in the continuous case, with $e^{2i\pi x}$, whose integral over an interval vanishes exactly on intervals of integer
length.

\writeup (a) Give the cell $(i,j)$, $0\le i<a$, $0\le j<b$, the weight $\zeta^{i+j}$. A horizontal brick covers cells
$(i,j),\dots,(i,j+k-1)$, of total weight $\zeta^{i+j}(1+\zeta+\dots+\zeta^{k-1})=0$, and likewise for a vertical brick. So the
total weight is $0$:
\[
0=\sum_{i<a}\sum_{j<b}\zeta^{i+j}=\Bigl(\sum_{i<a}\zeta^i\Bigr)\Bigl(\sum_{j<b}\zeta^j\Bigr).
\]
One factor vanishes. For $\zeta\ne1$, $\sum_{i<a}\zeta^i=\frac{\zeta^a-1}{\zeta-1}$ is $0$ if and only if $\zeta^a=1$, that is
$k\mid a$ (for $k=1$ there is nothing to prove). So $k\mid a$ or $k\mid b$. In particular a $10\times10$ board cannot be tiled
by $1\times4$ bricks, although $4$ divides its area.

(b) For a rectangle $R=[x_0,x_1]\times[y_0,y_1]$, Fubini for continuous functions (Tome~4) gives
\[
\iint_Re^{2i\pi(x+y)}\dd x\dd y=\int_{x_0}^{x_1}e^{2i\pi x}\dd x\int_{y_0}^{y_1}e^{2i\pi y}\dd y,
\]
and $\int_{x_0}^{x_1}e^{2i\pi x}\dd x=e^{2i\pi x_0}\frac{e^{2i\pi(x_1-x_0)}-1}{2i\pi}$ vanishes if and only if $x_1-x_0\in\Z$. So
the integral over $R$ is $0$ if and only if $R$ has a side of integer length. Each small rectangle gives $0$. The small
rectangles cover the big one and overlap only along edges, which have area zero, so the integral over the big rectangle is
the sum of the integrals over the small ones, which is $0$. Hence the big rectangle has an integer side.
\tools Tome~0, Chapter~2 (geometric sums of roots of unity); Tome~4, Chapter~4 (multiple integrals).
\leads De Bruijn (1969) used this to decide when an $a\times b\times c$ box can be filled with $1\times2\times4$ bricks.
Stan Wagon's paper on (b) gives fourteen proofs, from graph theory to the Euler characteristic, and the integral one is
the shortest.
\end{solution}

\begin{solution}{pb:prouhet}
\idea The signs $t_n$ are built digit by digit, so their generating function is a product with one factor per binary
digit, and each factor vanishes at $x=1$. Moments $\sum t_nn^k$ are derivatives at $1$. A zero of order $N$ kills all
derivatives of order below $N$.

\writeup Each $n<2^N$ is $\sum_{j<N}\varepsilon_j2^j$ with $\varepsilon_j\in\{0,1\}$, and $t_n=\prod_j(-1)^{\varepsilon_j}$. Expanding
the product, $G(x)=\sum_{n<2^N}t_nx^n=\prod_{j=0}^{N-1}(1-x^{2^j})$. Each factor vanishes at $1$, so $G=(1-x)^NH(x)$ with
$H(1)=\prod_j(1+1+\dots+1)=\prod_j2^j=2^{N(N-1)/2}$, since $\frac{1-x^{2^j}}{1-x}=1+x+\dots+x^{2^j-1}$.

Let $D=x\frac d{dx}$. Then $D^kG(x)=\sum_nt_nn^kx^n$, so the sum we want is $(D^kG)(1)$. By induction on $k$,
$D^kG=\sum_{i=0}^ka_{k,i}x^iG^{(i)}$ with integers $a_{k,i}$ and $a_{k,k}=1$: apply $D$ to $x^iG^{(i)}$ to get
$ix^iG^{(i)}+x^{i+1}G^{(i+1)}$. Since $G$ has a zero of order $N$ at $1$, $G^{(i)}(1)=0$ for $i<N$.

(a) For $k<N$, every term of $(D^kG)(1)$ contains some $G^{(i)}(1)$ with $i\le k<N$, so it is $0$.

(b) For $k=N$, only the term $i=N$ survives: the sum is $G^{(N)}(1)$. From $G=(-1)^N(x-1)^NH$ and Leibniz's rule,
$G^{(N)}(1)=(-1)^NN!\,H(1)=(-1)^NN!\,2^{N(N-1)/2}$.

(c) Take $N=4$. The numbers $n<16$ with $t_n=1$ (an even number of binary ones) are
$A=\{0,3,5,6,9,10,12,15\}$, and the others are $B=\{1,2,4,7,8,11,13,14\}$. By (a), $\sum_Ad^k=\sum_Bd^k$ for $k=0,1,2,3$:
the common values are $8$, $60$, $620$ and $7200$.
\tools Tome~3, Chapter~3 and Tome~4, Chapter~2 (generating functions); Tome~0, Chapter~3 (multiplicity of roots).
\leads This is Prouhet's solution (1851) of what is now called the Prouhet--Tarry--Escott problem: find two different sets
with equal power sums up to degree $k$. The Thue--Morse sequence uses $2^{k}$ numbers in each set; the best possible size is
$k+1$ numbers in each set, and such \emph{ideal} solutions are known for $k\le9$ and $k=11$, but not for $k=10$ or any $k\ge12$. The same sequence appears in fair sharing (``who chooses next''),
in the Rudin--Shapiro-like signs of Chapter~1, and in symbolic dynamics.
\end{solution}

\begin{solution}{pb:lightsout}
\idea Pressing commutes with rotating the circle, so the pressing map is a circulant matrix over $\mathbb F_2$, that is,
multiplication by a polynomial in the ring $\mathbb F_2[x]/(x^n-1)$. It is invertible exactly when that polynomial is prime
to $x^n-1$. For (b), a symmetric matrix has image orthogonal to its kernel, and the kernel is easy to guess.

\writeup Encode a set $u$ of lamps as $u(x)=\sum_{i\in u}x^i$ in $R=\mathbb F_2[x]/(x^n-1)$. Pressing the lamps of $u$
switches each lamp as many times as it has pressed lamps among itself and its neighbours, so it adds
$(x^{-1}+1+x)u=x^{-1}(1+x+x^2)u$ to the configuration (order does not matter, and pressing a lamp twice does nothing). As
$x$ is a unit in $R$, we must understand multiplication by $f=1+x+x^2$.

\emph{Lemma.} Over a field, multiplication by $f$ on $K[x]/(x^n-1)$ has a kernel of dimension $\deg g$, with
$g=\gcd(f,x^n-1)$. Indeed, $uf\equiv0$ means $x^n-1\mid uf$, that is $h\mid u\frac fg$ with $h=\frac{x^n-1}g$, and since
$\gcd(h,\frac fg)=1$ this means $h\mid u$. The multiples of $h$ of degree $<n$ are $hv$ with $\deg v<\deg g$.

(a) $1+x+x^2$ has no root in $\mathbb F_2$, so it is irreducible, and $g$ is $1$ or $1+x+x^2$. Since
$x^3-1=(x-1)(x^2+x+1)$, $x^3\equiv1$ modulo $f$, while $x\not\equiv1$ and $x^2\not\equiv1$. So $f\mid x^n-1$ if and only if
$3\mid n$. When $3\nmid n$, the kernel is $0$, pressing is a bijection of $R$, and every configuration can be switched off.
When $3\mid n$, the kernel has dimension $2$, so some configurations cannot.

(b) Let $3\mid n$, and let $v_r$ be the set of lamps $i\not\equiv r\pmod3$. Pressing $v_r$ switches each lamp twice or not at
all: a lamp in $v_r$ has exactly one neighbour in $v_r$, and a lamp outside it has both neighbours in it. So $v_0,v_1,v_2$
are in the kernel, which they span (two of them are independent and the kernel has dimension $2$). The pressing matrix
$A$ is symmetric, and over any field the image of $A$ is the orthogonal of the kernel of $A^T$ for the dot product (a
dimension count, Tome~0, Chapter~7, which holds over $\mathbb F_2$ as well). So a configuration $c$ can be switched off if
and only if $c\cdot v_r=0$ for $r=0,1,2$, that is $S_1+S_2$, $S_0+S_2$, $S_0+S_1$ are all even, that is $S_0\equiv S_1\equiv S_2
\pmod2$. The image has dimension $n-2$, so this is $\frac14$ of the $2^n$ configurations.
\tools Tome~2, Chapters~6 and~9 (polynomial rings, finite fields); Tome~0, Chapter~7 (rank and orthogonality).
\leads The commercial game \emph{Lights Out} (1995) is the $5\times5$ grid version. Its pressing matrix has a kernel of
dimension $2$ too, so only a quarter of the boards are solvable. Sutner (1989) showed that on every finite graph pressing
all vertices is enough to reach the all-off state from all-on: a theorem about $\mathbb F_2$-linear algebra.
\end{solution}

\begin{solution}{pb:sieving}
\idea At $q=1$ the $q$-binomial is the ordinary binomial, which counts all subsets. At a root of unity of order $m$, the
factors $1-q^i$ with $m\mid i$ vanish, and only they matter: they behave like $i$, divided by $m$. So the value at $\zeta$
looks like a binomial coefficient ``with everything divided by $m$''. On the other side, a subset fixed by a rotation of
order $m$ is a union of whole orbits of size $m$, and choosing $k/m$ orbits out of $n/m$ is the same binomial.

\writeup (a) $[n]_q=[k]_q+q^k[n-k]_q$. Divide by $[k]_q![n-k]_q!$ and multiply by $[n-1]_q!$:
\[
\qbinom nk=\frac{[n-1]_q!\,[k]_q}{[k]_q![n-k]_q!}+q^k\frac{[n-1]_q!\,[n-k]_q}{[k]_q![n-k]_q!}=\qbinom{n-1}{k-1}+q^k\qbinom{n-1}k.
\]
With $\qbinom n0=\qbinom nn=1$, induction on $n$ shows that $\qbinom nk$ is a polynomial with non-negative integer
coefficients.

(b) Cancel $[n-k]_q!$ and multiply each $[i]_q$ by $1-q$:
$\qbinom nk=\prod_{i=1}^k\frac{1-q^{n-k+i}}{1-q^i}$. The value of the polynomial at $\zeta$ is the limit of this product as
$q\to\zeta$. A factor $1-q^e$ vanishes at $\zeta$ if and only if $m\mid e$.

Suppose $m\mid k$. Then $m\mid n-k$, so $\zeta^{n-k+i}=\zeta^i$: the factors with $m\nmid i$ have the same non-zero limit above
and below, and cancel. The others are $i=mb$, $1\le b\le\frac km$; with $u=q^m\to1$,
\[
\frac{1-q^{n-k+mb}}{1-q^{mb}}=\frac{1-u^{(n-k)/m+b}}{1-u^b}\longrightarrow\frac{(n-k)/m+b}b,
\]
and $\prod_{b=1}^{k/m}\frac{(n-k)/m+b}b=\binom{n/m}{k/m}$.

Suppose $m\nmid k$. In the denominator, $m\mid i$ for $\lfloor k/m\rfloor$ values of $i\in\{1,\dots,k\}$. In the numerator,
$n-k+i$ runs over $\{n-k+1,\dots,n\}$, which contains $\lfloor\frac{k-1}m\rfloor+1$ multiples of $m$ (since $m\mid n$, they are
$n,n-m,\dots$ down to $n-k+1$). As $m\nmid k$, $\lfloor\frac{k-1}m\rfloor=\lfloor\frac km\rfloor$, so the numerator has one
more vanishing factor, of order one each. The limit is $0$.

(c) Let $d=\gcd(j,n)$. The subgroup generated by $g^j$ is generated by $g^d$, the rotation by $d$, whose orbits on
$\Z/n\Z$ are the $d$ classes modulo $d$, each of size $m=n/d$. A subset is fixed by $g^j$ if and only if it is fixed by
$g^d$, that is, it is a union of orbits. There are $\binom d{k/m}$ such subsets if $m\mid k$, and none otherwise. Since
$e^{2i\pi j/n}$ is a primitive $m$-th root of unity, this is the value given by (b).

(d) By Burnside's lemma (Tome~2, Chapter~3), the number of orbits is the average number of fixed points. With $F$ the polynomial $\qbinom nk$, it is
$\frac1n\sum_{j=0}^{n-1}F(e^{2i\pi j/n})=\frac1n\sum_{m\mid\gcd(n,k)}\varphi(m)\binom{n/m}{k/m}$. For $n=6$, $k=3$:
$\frac16\bigl(20+2\binom21\bigr)=4$.
\tools Tome~2, Chapter~3 (Burnside's lemma); Tome~0, Chapter~2 (roots of unity); Tome~4, Chapter~1 (finite sums).
\leads Reiner, Stanton and White (2004) named this the \emph{cyclic sieving phenomenon}: a polynomial that counts a set at
$q=1$ and counts fixed points of a cyclic action at roots of unity. Dozens of instances are known, for triangulations of
polygons, plane partitions and Young tableaux, and each one comes from representation theory.
\end{solution}

\begin{solution}{pb:phipm1}
\idea At $x=1$, the identity $1+x+\dots+x^{n-1}=\prod_{d\mid n,d>1}\Phi_d(x)$ becomes $n=\prod\Phi_d(1)$, a
multiplicative version of $\log n=\sum_{d\mid n}\Lambda(d)$: by Möbius, $\log\Phi_n(1)=\Lambda(n)$. At $x=-1$, two rules
move the index: substituting $-x$ doubles odd indices, and substituting $x^p$ multiplies by $p$ an index that $p$
already divides. Each rule is proved by comparing roots.

\writeup (a) For $n\ge2$, $\frac{x^n-1}{x-1}=\prod_{d\mid n,d>1}\Phi_d(x)$; at $x=1$, $n=\prod_{d\mid n,d>1}\Phi_d(1)$. All values
are positive (the left side is, and by induction each $\Phi_d(1)$ is). So $L(d)=\log\Phi_d(1)$ for $d\ge2$, $L(1)=0$,
satisfies $\sum_{d\mid n}L(d)=\log n=\sum_{d\mid n}\Lambda(d)$ for all $n\ge1$ (Tome~1, Chapter~4). By Möbius inversion, $L=\Lambda$:
$\Phi_n(1)=e^{\Lambda(n)}$.

(b) Two rules. \emph{(i) If $m>1$ is odd, $\Phi_{2m}(x)=\Phi_m(-x)$.} The roots of $\Phi_m(-x)$ are $-\zeta$ for $\zeta$ of order
$m$, and $-\zeta$ has order $2m$ ($m$ odd); there are $\varphi(m)=\varphi(2m)$ of them, distinct, and $\Phi_m(-x)$ is monic
since $\varphi(m)$ is even. Two monic polynomials of the same degree with the same simple roots are equal.
\emph{(ii) If $p\mid m$, $\Phi_{pm}(x)=\Phi_m(x^p)$.} If $z^p$ has order $m$, the order $e$ of $z$ satisfies
$e/\gcd(e,p)=m$, so $e=m$ or $pm$, and $e=m$ would need $\gcd(m,p)=1$. So the $p\varphi(m)$ distinct roots of $\Phi_m(x^p)$ have
order $pm$, and $\varphi(pm)=p\varphi(m)$ when $p\mid m$: same conclusion.

Now let $n\ge3$. If $n$ is odd, $\Phi_n(-1)=\Phi_{2n}(1)=1$ by (i) and (a). If $n=2m$ with $m\ge3$ odd,
$\Phi_n(-1)=\Phi_m(1)$, which is $p$ if $m$ is a power of $p$ and $1$ otherwise. If $4\mid n$, (ii) with $p=2$ gives
$\Phi_n(-1)=\Phi_{n/2}(1)$, which is $2$ if $n$ is a power of $2$ and $1$ otherwise. In short, for $n\ge3$:
$\Phi_n(-1)=p$ if $n=2p^k$ with $p$ an odd prime, $2$ if $n$ is a power of $2$, and $1$ otherwise.

(c) $\Phi_n(1)=\prod_{\gcd(a,n)=1}(1-\zeta^a)=1$ by (a). Each factor is in $\Z[\zeta]$, so $1-\zeta$ times the product of
the other factors is $1$.
\tools Tome~1, Chapters~4 and~5 (Möbius inversion, cyclotomic polynomials); Tome~0, Chapter~3 (polynomials and their roots).
\leads The units $\frac{1-\zeta^a}{1-\zeta}$ and the $1-\zeta$ of (c) are the \emph{cyclotomic units}. Kummer used them in
his attack on Fermat's last theorem, and they are linked to the values $L(1,\chi)$ by the class number formula (Tome~12).
When $n=p^k$, on the contrary, $1-\zeta$ generates the unique prime above $p$.
\end{solution}

\begin{solution}{pb:phipq}
\idea $\Phi_{pq}$ is a ratio of four binomials $1-x^d$. The denominator $\frac1{(1-x^p)(1-x^q)}$ is the generating function
of the numbers $ip+jq$, and below $pq$ each number has at most one such representation. Since $\Phi_{pq}$ has degree less
than $pq$, we lose nothing by working modulo $x^{pq}$, where $1-x^{pq}$ is $1$. For $\Phi_{105}$ the same reduction modulo
$x^8$ leaves a two-line computation.

\writeup From $x^n-1=\prod_{d\mid n}\Phi_d$ and Möbius inversion in multiplicative form,
$\Phi_n=\prod_{d\mid n}(x^d-1)^{\mu(n/d)}$, and for $n>1$, $\sum_{d\mid n}\mu(n/d)=0$, so the signs can be changed:
$\Phi_n=\prod_{d\mid n}(1-x^d)^{\mu(n/d)}$, an identity of rational functions, and of power series in $\Z[[x]]$.

(a) For $n=pq$: $\Phi_{pq}=\frac{(1-x)(1-x^{pq})}{(1-x^p)(1-x^q)}$. In $\Z[[x]]$,
$\frac1{(1-x^p)(1-x^q)}=\sum_{i,j\ge0}x^{ip+jq}$. If $ip+jq=i'p+j'q<pq$ then $p(i-i')=q(j'-j)$, so $q\mid i-i'$; as
$0\le i,i'<q$, $i=i'$ and then $j=j'$. So modulo $x^{pq}$,
\[
\Phi_{pq}\equiv(1-x)\sum_{m<pq}\mathbf 1_S(m)x^m,\qquad S=\{ip+jq:i,j\ge0\}.
\]
Since $\deg\Phi_{pq}=(p-1)(q-1)<pq$, the coefficient of $x^m$ in $\Phi_{pq}$ is exactly $\mathbf 1_S(m)-\mathbf 1_S(m-1)$,
for every $m<pq$. This is $-1$, $0$ or $1$. Moreover it is $1$ where $\mathbf 1_S$ jumps up and $-1$ where it jumps down,
and up and down jumps alternate: the non-zero coefficients alternate between $1$ and $-1$.

(b) $n=105=3\cdot5\cdot7$. Modulo $x^8$ only the divisors $d<8$ matter: $d=1,3,5,7$, with $\mu(105)=-1$, $\mu(35)=\mu(21)=\mu(15)
=1$. So
\[
\Phi_{105}\equiv\frac{(1-x^3)(1-x^5)(1-x^7)}{1-x}\equiv(1+x+\dots+x^7)(1-x^3-x^5-x^7)\pmod{x^8},
\]
since $x^8\equiv0$. The coefficient of $x^7$ is $1-1-1-1=-2$.
\tools Tome~1, Chapters~4 and~5 (Möbius inversion, cyclotomic polynomials); Tome~3, Chapter~3 and Tome~4, Chapter~2
(generating functions, the coin problem).
\leads The set $S$ is the set of amounts payable with coins $p$ and $q$; its complement has $\frac{(p-1)(q-1)}2$ elements, the
largest being $pq-p-q$ (Sylvester's coin problem, 1884), and this is the degree of $\Phi_{pq}$ minus the last jump. This
proof is due to Lam and Leung (1996). Schur's proof that the coefficients are unbounded uses $\Phi_{p_1\cdots p_t}$ with
$t$ odd and $p_1<\dots<p_t<p_1+p_2$: its coefficient of $x^{p_t}$ is $1-t$.
\end{solution}

\begin{solution}{pb:vanish}
\idea Relations between roots of unity are polynomials divisible by the minimal polynomial. For a prime $p$, the minimal
polynomial $\Phi_p$ has all coefficients $1$, which forces the coefficients to be equal. For $n=6$, a basis of
$\Q(\zeta_6)$ turns the question into two linear equations in $0$/$1$ unknowns. For $n=p^a$, $\zeta_{p^a}$ generates
$\Q(\zeta_{p^a})$ over $\Q(\omega)$ with a binomial minimal polynomial, so a relation splits into relations among $p$-th roots.

\writeup (a) Let $A=\sum_{k<p}a_kX^k\in\Q[X]$, of degree $<p$. If $A(\zeta)=0$, the minimal polynomial $\Phi_p$ of $\zeta$
(irreducible, Tome~0, Chapter~3) divides $A$; by degrees $A=c\Phi_p$, so all $a_k=c$. Conversely
$c(1+\zeta+\dots+\zeta^{p-1})=0$. For sets ($a_k\in\{0,1\}$) the only vanishing sums are the empty set and all $p$ roots.

(b) Let $\zeta=e^{i\pi/3}$, so $\zeta^3=-1$, and let $c_k\in\{0,1\}$ indicate whether $\zeta^k$ is in the set. Put
$u_k=c_k-c_{k+3}$ for $k=0,1,2$: the sum is $u_0+u_1\zeta+u_2\zeta^2$. As $\zeta^2=\zeta-1$, this is
$(u_0-u_2)+(u_1+u_2)\zeta$, and since $1,\zeta$ are linearly independent over $\R$, the sum vanishes if and only if
$u_0=u_2=-u_1=t$. If $t=0$, the set is a union of antipodal pairs $\{\zeta^k,-\zeta^k\}$: $7$ non-empty sets. If $t=1$,
$c_0=c_2=c_4=1$ and $c_1=c_3=c_5=0$: the triangle $\{1,\zeta^2,\zeta^4\}$. If $t=-1$, the other triangle. That makes $9$.

(c) Let $\zeta=e^{2i\pi/p^a}$, $N=p^{a-1}$ and $\omega=\zeta^N$, a primitive $p$-th root. By the tower law and
$[\Q(\zeta):\Q]=\varphi(p^a)$ (Tome~2, Chapter~10), $[\Q(\zeta):\Q(\omega)]=\frac{\varphi(p^a)}{p-1}=N$, and $\zeta$ is a root
of $X^N-\omega$, so that is its minimal polynomial over $\Q(\omega)$ and $1,\zeta,\dots,\zeta^{N-1}$ is a basis of $\Q(\zeta)$
over $\Q(\omega)$ (Tome~2, Chapter~8). Write each exponent as $r+Nj$, $0\le r<N$, $0\le j<p$. A set with indicator $c$ has
sum $\sum_{r<N}\zeta^r\bigl(\sum_{j<p}c_{r+Nj}\omega^j\bigr)$. It vanishes if and only if each inner sum vanishes, and by (a)
this means that for each $r$, either all $\zeta^{r+Nj}=\zeta^r\omega^j$ are in the set or none. These are the sets
$\{\eta,\eta\omega,\dots,\eta\omega^{p-1}\}$ with $\eta=\zeta^r$.
\tools Tome~0, Chapter~3 (irreducibility of $\Phi_p$); Tome~2, Chapters~8 and~10 (tower law, cyclotomic fields).
\leads For $n$ with two prime factors, (b) generalises: every vanishing sum of $n$-th roots with non-negative coefficients
is a sum of rotated regular $p$-gons and $q$-gons (de Bruijn, 1953), but with three primes, as for $\Phi_{105}$, this
fails, and cancellations between polygons appear. Lam and Leung (2000) proved that a vanishing sum of $m$ $n$-th roots of
unity, repetitions allowed, exists if and only if $m$ is a combination with non-negative integer coefficients of the
primes dividing $n$.
\end{solution}

\begin{solution}{pb:kronecker}
\idea Raising the roots to the $k$-th power keeps them in the closed unit disc and keeps the polynomial in $\Z[X]$, of the
same degree. But there are only finitely many such polynomials: their coefficients are bounded by Viète. So the powers
$\alpha^k$ take finitely many values, and $\alpha^k=\alpha^l$ for some $k<l$. To see that $\prod(X-\alpha_i^k)$ has integer
coefficients without symmetric functions, take a matrix.

\writeup (a) Let $d=\deg P$, with roots $\alpha_1,\dots,\alpha_d$ (with multiplicity), and $C\in M_d(\Z)$ its companion
matrix, whose characteristic polynomial is $P$. Over $\C$, $C$ is similar to an upper triangular matrix $T$ with diagonal
$\alpha_1,\dots,\alpha_d$ (Jordan form, Tome~2, Chapter~7). Then $C^k$ is similar to $T^k$, upper triangular with diagonal
$\alpha_i^k$, so $P_k=\prod_i(X-\alpha_i^k)$ is the characteristic polynomial of $C^k\in M_d(\Z)$: $P_k\in\Z[X]$, monic of degree
$d$. Its coefficients are, up to sign, the elementary symmetric functions of the $\alpha_i^k$ (Tome~0, Chapter~3), and since
$|\alpha_i^k|\le1$, the coefficient of $X^{d-j}$ has modulus at most $\binom dj$. There are finitely many monic integer
polynomials of degree $d$ with these bounds, so the set $R$ of all their roots is finite. For each $i$, the numbers
$\alpha_i,\alpha_i^2,\alpha_i^3,\dots$ all lie in $R$, so $\alpha_i^k=\alpha_i^l$ for some $k<l$; as $\alpha_i\ne0$ (because
$P(0)\ne0$), $\alpha_i^{l-k}=1$.

If $P$ is irreducible, its root $\alpha_1$ is a root of unity, of some order $n$, so it is a root of $\Phi_n\in\Z[X]$, which is
irreducible (Tome~2, Chapter~10). Two monic irreducible polynomials of $\Q[X]$ with a common root are equal: $P=\Phi_n$.

(b) Let $Q(x)=x^dP(x+\frac1x)$. Expanding $(x+\frac1x)^j=\sum_i\binom jix^{j-2i}$, $Q$ is a monic integer polynomial of
degree $2d$ with $Q(0)=1$. If $\beta=2\cos\theta$ is a root of $P$, then $x+\frac1x=\beta$ has the roots $e^{\pm i\theta}$, of
modulus $1$. Since $P=\prod(X-\beta_i)$, $Q=\prod_i(x^2-\beta_ix+1)$, and all its roots are of this form. By (a) they are
roots of unity: $e^{i\theta}=e^{2i\pi r}$ with $r\in\Q$, and $\beta=2\cos2\pi r$.

(c) $z+\bar z=\frac65$ and $z\bar z=1$, so $z$ is a root of $5X^2-6X+5$, irreducible over $\Q$ (negative discriminant).
If $z$ were a root of unity, it would be a root of some $\Phi_n$, so $5X^2-6X+5$ would divide $\Phi_n$ in $\Q[X]$; by Gauss's
lemma (Tome~0, Chapter~3) a primitive divisor of a monic integer polynomial is monic up to sign, and $5X^2-6X+5$ is
primitive and not monic. So $z$ is not a root of unity. What fails is integrality: the minimal polynomial is not in
$\Z[X]$ once made monic, $X^2-\frac65X+1$.
\tools Tome~0, Chapters~3 and~7 (Viète, Gauss's lemma, characteristic polynomials); Tome~2, Chapters~7 and~10 (Jordan form,
irreducibility of $\Phi_n$).
\leads Kronecker proved this in 1857. Lehmer asked in 1933 whether the \emph{Mahler measure} $\prod\max(1,|\alpha_i|)$ of a
non-cyclotomic irreducible integer polynomial (see Problem~\ref{pb:jensen}) is bounded away from $1$. The smallest known value,
$1.17628\ldots$, belongs to Lehmer's own polynomial $x^{10}+x^9-x^7-x^6-x^5-x^4-x^3+x+1$, and the question is still open
(Tome~12).
\end{solution}

\begin{solution}{pb:primsum}
\idea Inside $\mathbb F_p$, a primitive root $\gamma$ plays the role of $e^{2i\pi/(p-1)}$: it has order exactly $p-1$, its
powers are all of $\mathbb F_p^*$, and the primitive roots are the $\gamma^a$ with $a$ prime to $p-1$. So the sum of their
$k$-th powers is a Ramanujan sum computed in $\mathbb F_p$, and the proof of Problem~\ref{pb:ramsum}(a) works word for word,
because it only used the geometric sum.

\writeup Let $q=p-1$ and let $\gamma$ be a primitive root (Tome~1, Chapter~5). The primitive roots modulo $p$ are the
$\gamma^a$, $1\le a\le q$, $\gcd(a,q)=1$, so in $\mathbb F_p$ the sum of their $k$-th powers is
$C_q(k)=\sum_{\gcd(a,q)=1}\gamma^{ak}$. For $e\mid q$ put $\gamma_e=\gamma^{q/e}$, an element of order $e$, and
$C_e(k)=\sum_{1\le b\le e,\ \gcd(b,e)=1}\gamma_e^{bk}$. The geometric sum $\sum_{a=1}^e\gamma_e^{ak}$ is $e$ if $e\mid k$ (all
terms are $1$) and $\frac{\gamma_e^{k}(\gamma_e^{ek}-1)}{\gamma_e^k-1}=0$ otherwise. Sorting $a$ by $\gcd(a,e)$ as in
Problem~\ref{pb:ramsum}(a), and noting that $\gamma_e^{e/f}=\gamma_f$, gives $\sum_{f\mid e}C_f(k)=e\cdot\mathbf 1_{e\mid k}$ in
$\mathbb F_p$, for all $e\mid q$. Möbius inversion over the divisors of $q$ (valid in any abelian group) gives
$C_q(k)=\sum_{d\mid\gcd(q,k)}\mu(q/d)d=c_q(k)$, which is (c).

(a) is $k=1$: $c_q(1)=\mu(q)$. For (b), $g\mapsto g^{-1}$ permutes the primitive roots, and $g\ne g^{-1}$ because $g^2=1$
would mean $q\le2$, that is $p=3$. So the primitive roots split into pairs $\{g,g^{-1}\}$ with product $1$.
\tools Tome~1, Chapters~4 and~5 (Möbius inversion, primitive roots); Tome~2, Chapter~9 (cyclic multiplicative group).
\leads Gauss stated (a) in the \emph{Disquisitiones} (article~81). Since $\mu(p-1)=0$ as soon as $4\mid p-1$, the sum of the
primitive roots is then divisible by $p$. Counting primitive roots in short intervals, where these sums meet character sums,
is still an active field (Tome~5).
\end{solution}

\begin{solution}{pb:phimodp}
\idea Modulo $p$, the polynomial $X^n-1$ still factors as $\prod_{d\mid n}\Phi_d$, and if $p\nmid n$ it has no repeated root.
So the roots of $\Phi_n$ in a big enough finite field are exactly the elements of order $n$. The degree of such an element
over $\mathbb F_p$ is the size of the smallest field $\mathbb F_{p^f}$ that contains it, and $\mathbb F_{p^f}^*$ is cyclic of order
$p^f-1$. So everything is decided by when $n$ divides $p^f-1$.

\writeup (a) Reduce $X^n-1=\prod_{d\mid n}\Phi_d$ modulo $p$. Since $p\nmid n$, the derivative $nX^{n-1}$ is prime to $X^n-1$, so
$X^n-1$ has distinct roots in a splitting field $\mathbb F$ (Tome~2, Chapter~9). A root $\zeta$ of $\Phi_n$ has order dividing
$n$; if its order were $d<n$, $\zeta$ would also be a root of $X^d-1=\prod_{e\mid d}\Phi_e$, hence of some $\Phi_e$ with $e\ne n$,
and a double root of $X^n-1$. So the $\varphi(n)$ roots of $\Phi_n$ modulo $p$ are distinct and of order exactly $n$. For such
a $\zeta$, the degree of its minimal polynomial over $\mathbb F_p$ is $[\mathbb F_p(\zeta):\mathbb F_p]$, the least $f$ with
$\zeta\in\mathbb F_{p^f}$, that is with $\zeta^{p^f-1}=1$, that is $n\mid p^f-1$: $f=\operatorname{ord}_n(p)$. So every irreducible
factor of $\Phi_n$ modulo $p$ has degree $f$, they are distinct (no repeated roots), and there are $\varphi(n)/f$ of them.

(b) $X^4+1=\Phi_8$ is irreducible over $\Q$ (Tome~2, Chapter~10, or Eisenstein after $X\mapsto X+1$). Modulo $2$,
$X^4+1=(X+1)^4$. For odd $p$, $p^2\equiv1\pmod8$, so $f=\operatorname{ord}_8(p)\le2$, and $\Phi_8$ splits into at least two factors.

(c) First let $a=1$. Since $p\nmid m$, $\Phi_m(X^p)=\Phi_{pm}(X)\Phi_m(X)$: both sides are monic of degree $p\varphi(m)$, and
the roots of the left side, the $z$ with $z^p$ of order $m$, have order $m$ or $pm$, all simple. Modulo $p$,
$\Phi_m(X^p)\equiv\Phi_m(X)^p$ (Frobenius: $(\sum c_iX^i)^p\equiv\sum c_iX^{ip}$ for $c_i\in\mathbb F_p$). Cancelling $\Phi_m$ in the
domain $\mathbb F_p[X]$, $\Phi_{pm}\equiv\Phi_m^{p-1}$. For $a\ge2$, rule (ii) of Problem~\ref{pb:phipm1} gives
$\Phi_{p^am}(X)=\Phi_{pm}(X^{p^{a-1}})\equiv\Phi_{pm}(X)^{p^{a-1}}\equiv\Phi_m^{(p-1)p^{a-1}}$.
\tools Tome~2, Chapters~9 and~10 (finite fields, cyclotomic polynomials); Tome~1, Chapter~5 (orders).
\leads (a) is the simplest case of the decomposition of primes in number fields: $p$ splits in $\Q(\zeta_n)$ into
$\varphi(n)/f$ primes of degree $f$ (Tome~12). $X^4+1$ shows that irreducibility over $\Q$ cannot always be detected modulo a
single prime, a fact that every polynomial factoring algorithm must live with.
\end{solution}

\begin{solution}{pb:rattan}
\idea A root of unity in a field drags the full cyclotomic field of the lcm into it, and degrees ($\varphi$) cannot grow. The
field $\Q(i)=\Q(\zeta_4)$ thus contains only the four fourth roots of unity, and a rational tangent or a lattice polygon
produces a root of unity in $\Q(i)$. For cosines, use integrality instead of fields.

\writeup (a) Let $\zeta_m\in K=\Q(\zeta_n)$ and $L=\operatorname{lcm}(m,n)$, $g=\gcd(m,n)$. Bézout gives $um+vn=g$, and
$\zeta_m^v\zeta_n^u=e^{2i\pi(vn+um)/mn}=e^{2i\pi g/mn}=\zeta_L$. So $\Q(\zeta_L)\subseteq K$, and comparing degrees (Tome~2,
Chapter~10) $\varphi(L)\le\varphi(n)$. Write $L=nt$. The formula $\varphi(N)=N\prod_{p\mid N}(1-\frac1p)$ gives
\[
\varphi(nt)=\varphi(n)\,t\prod_{p\mid t,\ p\nmid n}\Bigl(1-\frac1p\Bigr)\ge\varphi(n)\,\varphi(t).
\]
So $\varphi(t)\le1$, $t\in\{1,2\}$, and $t=2$ needs $2\cdot(1-\frac12)=1$, that is $n$ odd. Hence $m\mid L$, with $L=n$ or ($L=2n$
and $n$ odd). Conversely, $-1\in K$, so $K$ contains the $2n$-th roots of unity $\pm\zeta_n^j$.

(b) Let $t=\tan\theta$ ($\cos\theta\ne0$). Then $\frac{1+it}{1-it}=\frac{\cos\theta+i\sin\theta}{\cos\theta-i\sin\theta}=e^{2i\theta}$,
a root of unity (as $\theta/\pi\in\Q$) in $\Q(i)=\Q(\zeta_4)$. By (a) it is a fourth root of unity, so
$2\theta\in\frac\pi2\Z$, $\theta\in\frac\pi4\Z$, and $\tan\theta\in\{0,\pm1\}$.

(c) Let $A,B,C$ be consecutive vertices of a regular $n$-gon, as Gaussian integers. Then $C-B$ is $B-A$ rotated by
$\pm\frac{2\pi}n$, so $\frac{C-B}{B-A}=e^{\pm2i\pi/n}\in\Q(i)$, a root of unity of order $n$. By (a), $n\mid4$, so $n=4$. The
square with vertices $0,1,1+i,i$ exists.

(d) $2\cos\theta=\zeta+\zeta^{-1}$ with $\zeta=e^{i\theta}$ a root of unity. Roots of unity are algebraic integers, and so are
their sums (Tome~2, Chapter~11). A rational algebraic integer is an integer, so $2\cos\theta\in\{-2,-1,0,1,2\}$.
\tools Tome~2, Chapters~10 and~11 (cyclotomic fields, algebraic integers); Tome~1, Chapters~1 and~4 (Bézout, $\varphi$).
\leads (d) is Niven's theorem (1956). In $\Z^3$, there are regular triangles and hexagons (in the plane $x+y+z=0$) as well as
squares, but still no regular pentagon: the angle $\frac{2\pi}5$ would need $\cos\frac{2\pi}5\in\Q$. The same reasoning is
the crystallographic restriction of Chapter~1.
\end{solution}

\begin{solution}{pb:ramacube}
\idea Call the cube roots $a,b,c$. We know their cubes through the symmetric functions of the cosines. The identity
$s^3=a^3+b^3+c^3+3st-3abc$, with $s=a+b+c$ and $t=ab+bc+ca$, gives $s^3$ in terms of $st$; the same identity for $ab,bc,ca$
gives $t^3$. Multiplying, $m=st$ satisfies a cubic, which turns out to be a pure cube after a shift.

\writeup By Tome~2, Chapter~10 (or directly from $1+\zeta+\dots+\zeta^6=0$), $2\cos\frac{2\pi}7$, $2\cos\frac{4\pi}7$,
$2\cos\frac{8\pi}7$ are the roots of $u^3+u^2-2u-1$; so the three cosines $x_1,x_2,x_3$ are the roots of $8x^3+4x^2-4x-1$:
\[
x_1+x_2+x_3=-\tfrac12,\qquad x_1x_2+x_2x_3+x_3x_1=-\tfrac12,\qquad x_1x_2x_3=\tfrac18.
\]
Let $a,b,c$ be the real cube roots of $x_1,x_2,x_3$, so $abc=\frac12$; let $s=a+b+c$, $t=ab+bc+ca$, real. For any $a,b,c$,
$(a+b+c)^3=a^3+b^3+c^3+3(a+b+c)(ab+bc+ca)-3abc$. Hence
\[
s^3=-\tfrac12+3st-\tfrac32=3st-2.
\]
The same identity for $ab,bc,ca$, whose sum is $t$, whose pairwise products sum to $abc\cdot s$, and whose product is
$(abc)^2$, gives
\[
t^3=(x_1x_2+x_2x_3+x_3x_1)+3t\cdot\tfrac12s-3\cdot\tfrac14=\tfrac32st-\tfrac54.
\]
Let $m=st$. Multiplying, $m^3=(3m-2)(\frac32m-\frac54)$, that is $4m^3-18m^2+27m-10=0$. Put $m=y+\frac32$: the terms in $y^2$
and $y$ cancel, and the equation becomes $4y^3+\frac72=0$, $y^3=-\frac78$. As $m$ is real, $y=-\frac{\sqrt[3]7}2$ and
$m=\frac{3-\sqrt[3]7}2$. Then $s^3=3m-2=\frac{5-3\sqrt[3]7}2$, and since $s$ is real, $s=\sqrt[3]{\frac{5-3\sqrt[3]7}2}$.
Numerically both sides are $-0.71751\ldots$
\tools Tome~0, Chapter~3 (Viète, cubic equations); Tome~2, Chapter~10 (the period polynomial of the heptagon).
\leads Ramanujan wrote this identity, and the companion
$\sqrt[3]{\cos\frac{2\pi}9}+\sqrt[3]{\cos\frac{4\pi}9}+\sqrt[3]{\cos\frac{8\pi}9}=\sqrt[3]{\frac{3\sqrt[3]9-6}2}$, in his notebooks
(Berndt, \emph{Ramanujan's Notebooks}, Part~IV). The method works for the periods of any cyclic cubic field, and explains
why such ``denesting'' miracles exist: the cubic for $m$ has a cyclic Galois group, so its roots are expressible by real
radicals of a special shape.
\end{solution}

\begin{solution}{pb:discphi}
\idea The discriminant is a product over pairs of roots, which is the product of the derivatives at the roots, and the
derivative of $\Phi_p$ at a root is easy from $(x-1)\Phi_p=x^p-1$. Its square root is a product of differences of roots of
unity, so it lies in $\Q(\zeta)$. For Zolotarev, the sign of $\pi_a$ is a character of a cyclic group, so it is determined
by its value at a generator, which is the sign of a $(p-1)$-cycle. And the automorphism $\sigma_a$ permutes the factors
of $\prod_{i<j}(\zeta^j-\zeta^i)$ exactly as $\pi_a$ permutes $\{1,\dots,p-1\}$.

\writeup (a) For a monic $P$ of degree $d$ with simple roots $\alpha_i$, $P'(\alpha_i)=\prod_{j\ne i}(\alpha_i-\alpha_j)$, so
\[
\prod_iP'(\alpha_i)=\prod_{i\ne j}(\alpha_i-\alpha_j)=(-1)^{d(d-1)/2}\prod_{i<j}(\alpha_i-\alpha_j)^2.
\]
For $P=\Phi_p$,
$d=p-1$ and $\frac{(p-1)(p-2)}2\equiv\frac{p-1}2\pmod2$ since $p-2$ is odd. Differentiating $(x-1)\Phi_p(x)=x^p-1$ at a root
$\zeta^i$ gives $(\zeta^i-1)\Phi_p'(\zeta^i)=p\zeta^{i(p-1)}=p\zeta^{-i}$. Over the $p-1$ roots,
$\prod_i\zeta^{-i}=1$ (the product of all $p$-th roots of unity is $1$ for odd $p$) and
$\prod_i(\zeta^i-1)=\prod_i(1-\zeta^i)=\Phi_p(1)=p$ ($p-1$ is even). So $\prod_i\Phi_p'(\zeta^i)=p^{p-1}/p=p^{p-2}$, and
\[
\operatorname{disc}\Phi_p=(-1)^{(p-1)/2}p^{p-2}=p^*\,p^{p-3}.
\]
Let $\delta=\prod_{1\le i<j\le p-1}(\zeta^j-\zeta^i)\in\Q(\zeta)$. Then $\delta^2=p^*p^{p-3}$, and since $p-3$ is even,
$\bigl(\delta/p^{(p-3)/2}\bigr)^2=p^*$: $\sqrt{p^*}=\pm\delta/p^{(p-3)/2}\in\Q(\zeta)$.

(b) $\pi_{ab}=\pi_a\circ\pi_b$, so $a\mapsto\operatorname{sgn}\pi_a$ is a homomorphism $(\Z/p\Z)^*\to\{\pm1\}$ (Tome~2, Chapter~4).
Let $g$ be a primitive root. $\pi_g$ fixes $0$ and permutes $1,g,g^2,\dots,g^{p-2}$ cyclically: a $(p-1)$-cycle, of sign
$(-1)^{p-2}=-1$. So $\operatorname{sgn}\pi_{g^k}=(-1)^k$, which is $1$ exactly when $g^k$ is a square, that is when $k$ is
even (Tome~1, Chapter~6). This is the Legendre symbol.

(c) $\sigma_a(\delta)=\prod_{i<j}(\zeta^{aj}-\zeta^{ai})$. As $(i,j)$ runs over pairs $i<j$ in $\{1,\dots,p-1\}$, the residues
$(ai,aj)$ run over all pairs of distinct non-zero residues, each unordered pair once, but in the wrong order exactly when
$\pi_a$ reverses it. So $\sigma_a(\delta)=(-1)^{\mathrm{inv}(\pi_a)}\delta=\operatorname{sgn}(\pi_a)\delta$, where $\mathrm{inv}$
counts inversions of $\pi_a$ on $\{1,\dots,p-1\}$ (its sign is that of $\pi_a$ on $\Z/p\Z$, as $0$ is fixed). With (b) and
(a), $\sigma_a(\sqrt{p^*})=\bigl(\frac ap\bigr)\sqrt{p^*}$.
\tools Tome~0, Chapter~3 (roots and derivatives); Tome~2, Chapters~4 and~10 (signs of permutations, Galois group of
$\Q(\zeta_p)$); Tome~1, Chapters~5 and~6 (primitive roots, Legendre symbol).
\leads (c) says that the quadratic subfield $\Q(\sqrt{p^*})$ of $\Q(\zeta_p)$ is the fixed field of the squares, and that
$\sigma_a$ acts on it through $\bigl(\frac ap\bigr)$. Taking $a=q$ another prime, and using that $\sigma_q$ is the Frobenius
at $q$, gives yet another proof of quadratic reciprocity, the one that generalises to Artin reciprocity (Tome~12). Zolotarev
(1872) used (b) to give a proof by permutations alone.
\end{solution}

\begin{solution}{pb:eisensine}
\idea Gauss's lemma counts the sign changes in $j\mapsto qj\bmod p$; the odd function $\sin\frac{2\pi x}p$ records the
signs as a product of sines. The polynomial identity (a) then turns each ratio $\frac{\sin qx}{\sin x}$ into a product of
differences $\sin^2\frac{2\pi j}p-\sin^2\frac{2\pi k}q$. The double product is symmetric in $p$ and $q$ except for the sign
of each factor, and counting the factors gives reciprocity.

\writeup (a) By de Moivre, $\sin qx=\operatorname{Im}(\cos x+i\sin x)^q=\sum_{j\text{ odd}}\binom qj(-1)^{(j-1)/2}\cos^{q-j}x\sin^jx$.
For odd $j$, $q-j$ is even and $\cos^{q-j}x=(1-\sin^2x)^{(q-j)/2}$. So $\frac{\sin qx}{\sin x}=F(\sin^2x)$ for a polynomial $F$
of degree $\frac{q-1}2$, with leading coefficient
$\sum_{j\text{ odd}}\binom qj(-1)^{(j-1)/2}(-1)^{(q-j)/2}=(-1)^{(q-1)/2}\sum_{j\text{ odd}}\binom qj=(-1)^{(q-1)/2}2^{q-1}$,
which is $(-4)^{(q-1)/2}$. For $1\le k\le\frac{q-1}2$, $x_k=\frac{2\pi k}q\in\left]0,\pi\right[$ satisfies $\sin qx_k=0\ne\sin x_k$, so
$\sin^2x_k$ is a root of $F$. These $\frac{q-1}2$ values are distinct: $\sin^2x_k=\sin^2x_l$ with $x_k,x_l\in\left]0,\pi\right[$ means
$x_l=x_k$ or $x_l=\pi-x_k$, and the second gives $2(k+l)=q$, impossible. So $F(y)=(-4)^{(q-1)/2}\prod_k(y-\sin^2x_k)$.

(b) Let $H=\{1,\dots,\frac{p-1}2\}$. For $j\in H$, write $qj\equiv\varepsilon_jr_j\pmod p$ with $r_j\in H$ and $\varepsilon_j=\pm1$.
The map $j\mapsto r_j$ is injective, since $r_j=r_{j'}$ gives $qj\equiv\pm qj'$, $j\equiv\pm j'$, and $0<j+j'<p$ leaves
only $j=j'$; so it is a bijection of $H$. Multiplying the congruences, $q^{(p-1)/2}\bigl(\frac{p-1}2\bigr)!\equiv
\prod_j\varepsilon_j\bigl(\frac{p-1}2\bigr)!$, so $q^{(p-1)/2}\equiv\prod_j\varepsilon_j$, and by Euler's criterion (Tome~1, Chapter~6)
$\bigl(\frac qp\bigr)=\prod_j\varepsilon_j$. On the other side, $x\mapsto\sin\frac{2\pi x}p$ is odd and $p$-periodic, so
$\sin\frac{2\pi qj}p=\varepsilon_j\sin\frac{2\pi r_j}p$, and
$\prod_j\frac{\sin(2\pi qj/p)}{\sin(2\pi j/p)}=\prod_j\varepsilon_j\cdot\frac{\prod_j\sin(2\pi r_j/p)}{\prod_j\sin(2\pi j/p)}
=\prod_j\varepsilon_j$, the sines being non-zero.

(c) Let $u_j=\sin^2\frac{2\pi j}p$ and $v_k=\sin^2\frac{2\pi k}q$, for $j\le\frac{p-1}2$, $k\le\frac{q-1}2$. By (b), then (a) with
$x=\frac{2\pi j}p$,
\[
\Bigl(\frac qp\Bigr)=\prod_j\frac{\sin q\frac{2\pi j}p}{\sin\frac{2\pi j}p}=(-4)^{\frac{p-1}2\frac{q-1}2}\prod_{j,k}(u_j-v_k).
\]
Exchanging the roles, $\bigl(\frac pq\bigr)=(-4)^{\frac{p-1}2\frac{q-1}2}\prod_{j,k}(v_k-u_j)$. The products differ by the sign
$(-1)^{\frac{p-1}2\frac{q-1}2}$, one for each factor, and are non-zero: $u_j=v_k$ would give $\frac{2\pi j}p=\frac{2\pi k}q$ or
$\frac{2\pi j}p=\pi-\frac{2\pi k}q$, that is $qj=pk$ (impossible, $p\nmid j$) or $2(qj+pk)=pq$ (impossible, $pq$ is odd).
Dividing the two equalities, $\bigl(\frac qp\bigr)\big/\bigl(\frac pq\bigr)=(-1)^{\frac{p-1}2\frac{q-1}2}$, and as both symbols
are $\pm1$, this is the law of reciprocity.
\tools Tome~0, Chapters~2 and~3 (de Moivre, polynomials and roots); Tome~1, Chapter~6 (Euler's criterion, Gauss's lemma).
\leads Eisenstein published this in 1845. The same pattern, with $\sin$ replaced by elliptic functions with complex
multiplication by $\Z[i]$ or $\Z[\omega]$, gave him the biquadratic and cubic reciprocity laws: the special functions do the
counting, and their factorisation does the rest. In modern language, (a) is a statement about the cyclotomic units of
Problem~\ref{pb:phipm1}.
\end{solution}

\begin{solution}{pb:jensen}
\idea Each quadratic factor is $|x-e^{i\varphi}|^2$, and the angles $\theta+\frac{2\pi k}n$ are the arguments of the $n$-th
roots of $e^{in\theta}$, so the product is $|x^n-e^{in\theta}|^2$. Taking logarithms gives a Riemann sum for Poisson's
integral, exactly computable. Jensen's formula is then the sum of Poisson's integrals over the roots.

\writeup (a) $x^2-2x\cos\varphi+1=(x-e^{i\varphi})(x-e^{-i\varphi})$, and the $e^{i(\theta+2\pi k/n)}$, $0\le k<n$, are the $n$
roots of $X^n=e^{in\theta}$. So $\prod_k(x-e^{i(\theta+2\pi k/n)})=x^n-e^{in\theta}$, and likewise with $-\theta$. The product
is $(x^n-e^{in\theta})(x^n-e^{-in\theta})=x^{2n}-2x^n\cos n\theta+1$.

(b) For $|x|\ne1$, $f(\theta)=\log(1-2x\cos\theta+x^2)$ is continuous (as $1-2x\cos\theta+x^2=|x-e^{i\theta}|^2>0$), so its
Riemann sums converge (Tome~0, Chapter~6). By (a) with $\theta=0$,
\[
\frac{2\pi}n\sum_{k=0}^{n-1}f\Bigl(\frac{2\pi k}n\Bigr)=\frac{2\pi}n\log(x^n-1)^2=\frac{4\pi}n\log|x^n-1|.
\]
If $|x|<1$, $\log|x^n-1|\to0$ and the limit is $0$. If $|x|>1$, $\frac{4\pi}n\log|x^n-1|=4\pi\log|x|+\frac{4\pi}n\log|1-x^{-n}|\to
4\pi\log|x|$. So $\int_0^{2\pi}f$ is $0$, or $4\pi\log|x|$, and as $f(2\pi-\theta)=f(\theta)$, $\int_0^\pi f$ is half of that.

(c) Write $z_k=r_ke^{i\alpha_k}$, $r_k\ne1$. Then $\log|e^{i\theta}-z_k|=\frac12\log(1-2r_k\cos(\theta-\alpha_k)+r_k^2)$. Over a
full period the shift by $\alpha_k$ does not matter, so by (b) (and its full-period version)
$\frac1{2\pi}\int_0^{2\pi}\log|e^{i\theta}-z_k|\dd\theta=\frac1{2\pi}\cdot\frac12\cdot4\pi\log\max(1,r_k)=\log\max(1,|z_k|)$. Add
these up, with $\log|a|$ for the constant.

The quantity $M(P)=|a|\prod_k\max(1,|z_k|)$ is the Mahler measure, and (c) says $\log M(P)$ is the average of
$\log|P|$ on the circle. If $P\in\Z[X]$ is irreducible with $P(0)\ne0$ and $M(P)=1$, then $|a|\ge1$ forces $|a|=1$ and all
$|z_k|\le1$: by Problem~\ref{pb:kronecker}, $P=\pm\Phi_n$.
\tools Tome~0, Chapters~2 and~6 (roots of unity, Riemann sums); Tome~4, Chapter~4 (integrals); Problem~\ref{pb:kronecker}.
\leads (c) is Jensen's formula (1899) on the unit circle, for polynomials. For holomorphic functions it counts zeros in a
disc, and Tome~5, Chapter~6 uses it to bound the zeros of $\zeta$. The formula also holds when $P$ has roots on the circle,
since the logarithmic singularities are integrable.
\end{solution}

\begin{solution}{pb:midpoint}
\idea Taking midpoints treats all vertices alike, so it is diagonal in the Fourier basis of $\C^n$: the mode
$k\mapsto\omega^{jk}$ is multiplied by $\frac{1+\omega^j}2=\cos\frac{\pi j}n\,e^{i\pi j/n}$. The centroid is mode $0$. The
largest remaining moduli are those of $j=1$ and $j=n-1$, both $\cos\frac\pi n$, so after many steps only these two modes are
visible, and two opposite modes make an affine image of a regular polygon.

\writeup (a) By discrete Fourier inversion (Tome~4, Chapter~7), $z_k=\sum_{j=0}^{n-1}\hat z_j\omega^{jk}$ with
$\hat z_j=\frac1n\sum_kz_k\omega^{-jk}$, and $\hat z_0=0$ is the centroid. For the mode $u_k=\omega^{jk}$,
$(Mu)_k=\frac12(1+\omega^j)\omega^{jk}$, and
$\lambda_j=\frac{1+\omega^j}2=e^{i\pi j/n}\frac{e^{-i\pi j/n}+e^{i\pi j/n}}2=\cos\frac{\pi j}n\,e^{i\pi j/n}$. So
$(M^tz)_k=\sum_j\lambda_j^t\hat z_j\omega^{jk}$. Now $\lambda_1=\cos\frac\pi n\,e^{i\pi/n}$ and
$\lambda_{n-1}=\cos\frac{\pi(n-1)}n\,e^{i\pi(n-1)/n}=\cos\frac\pi n\,e^{-i\pi/n}$ (both signs change). For $2\le j\le n-2$,
$|\lambda_j|=|\cos\frac{\pi j}n|<\cos\frac\pi n$. Dividing by $\cos^t\frac\pi n$, those modes tend to $0$, and what remains is
$a e^{i\pi t/n}\omega^k+be^{-i\pi t/n}\omega^{-k}$ with $a=\hat z_1$, $b=\hat z_{n-1}$.

(b) $L(u)=au+b\bar u$ is $\R$-linear on $\C=\R^2$, and $L(\omega^k)=a\omega^k+b\omega^{-k}$. On the basis $1,i$, writing
$a=a_1+ia_2$, $b=b_1+ib_2$, $L(1)=(a_1+b_1,a_2+b_2)$ and $L(i)=(b_2-a_2,a_1-b_1)$, so $\det L=(a_1+b_1)(a_1-b_1)-(a_2+b_2)(b_2-a_2)
=|a|^2-|b|^2$. The rescaled polygons approach images of the regular polygon (rotated by $\frac{\pi t}n$, which does not change
its shape) by an invertible linear map when $|a|\ne|b|$: polygons inscribed in an ellipse, at equally spaced parameters.

(c) For $n=4$, $\lambda_2=\frac{1+\omega^2}2=0$: after one step only the modes $1$ and $3$ remain, and $Mz$ is an affine image of
a square, a parallelogram. This is Varignon's theorem: the midpoints of any quadrilateral form a parallelogram.
\tools Tome~4, Chapter~7 (discrete Fourier transform); Tome~0, Chapters~2 and~7 (roots of unity, determinants).
\leads Darboux (1878) studied this iteration, and Berlekamp, Gilbert and Sinden (1965) used it for random polygons. The same
eigenvalue computation shows that ``averaging with neighbours'' on a cycle mixes like $\cos^t\frac\pi n$, a first example of
the spectral gap of a graph, and of diffusion on the circle (Tome~3, Chapter~8).
\end{solution}

\begin{solution}{pb:napoleon}
\idea This is the case $n=3$ of Problem~\ref{pb:pdn}, and it is worth doing by hand. The centres are apexes of isosceles
triangles with apex angle $120^\circ$, a circulant operation on the triangle. It kills one of the two non-trivial Fourier
modes of $\Z/3\Z$, and one mode alone is an equilateral triangle. For areas, the signed area of a polygon is a sum of
contributions $\pm|\hat z_j|^2$ of its modes.

\writeup Let $\omega=e^{2i\pi/3}$ and $z_0,z_1,z_2$ be $A,B,C$, counterclockwise. The point $p$ with
$z_k-p=\omega(z_{k+1}-p)$ is the centre of rotation by $120^\circ$ mapping $z_{k+1}$ to $z_k$, so it is the centre of the
equilateral triangle on $z_kz_{k+1}$; for $z_k=0$, $z_{k+1}=1$ it is $p=\frac{-\omega}{1-\omega}=\frac12-\frac i{2\sqrt3}$, below
the side, that is outside the counterclockwise triangle. So $p_k=\frac{z_k-\omega z_{k+1}}{1-\omega}$. Write
$z_k=\hat z_0+\hat z_1\omega^k+\hat z_2\omega^{2k}$. The mode $\omega^{jk}$ is multiplied by $\frac{1-\omega^{j+1}}{1-\omega}$:
$1$ for $j=0$, $\frac{1-\omega^2}{1-\omega}=1+\omega$ for $j=1$, and $0$ for $j=2$. So
\[
p_k=\hat z_0+(1+\omega)\hat z_1\,\omega^k,
\]
an equilateral triangle centred at $\hat z_0$, the centroid of $ABC$. Inwards, $\omega$ becomes $\omega^2$:
$p'_k=\frac{z_k-\omega^2z_{k+1}}{1-\omega^2}$ kills the mode $j=1$, and $p'_k=\hat z_0+c\,\hat z_2\omega^{2k}$ with
$c=\frac{1-\omega^4}{1-\omega^2}=1+\omega^2$, again equilateral with the same centre. This proves (a) and the first half of (b).

Areas. The signed area of a polygon is $\frac12\operatorname{Im}\sum_k\bar z_kz_{k+1}$ (the shoelace formula). With
$z_k=\sum_j\hat z_j\omega^{jk}$ and $\sum_k\omega^{(l-j)k}=3\cdot\mathbf 1_{j=l}$, we get $\sum_k\bar z_kz_{k+1}=3\sum_j|\hat z_j|^2\omega^j$,
so the signed area is $\frac{3\sqrt3}4\bigl(|\hat z_1|^2-|\hat z_2|^2\bigr)$. For $ABC$ this is its (positive) area. The outer
triangle has only the mode $1$, with $|1+\omega|=1$: area $\frac{3\sqrt3}4|\hat z_1|^2$. The inner one has only the mode $2$,
with $|1+\omega^2|=1$: signed area $-\frac{3\sqrt3}4|\hat z_2|^2$, so it is clockwise, of area $\frac{3\sqrt3}4|\hat z_2|^2$. The
difference of the areas is $\frac{3\sqrt3}4(|\hat z_1|^2-|\hat z_2|^2)$, the area of $ABC$.
\tools Tome~0, Chapter~2 (complex numbers and rotations); Tome~4, Chapter~7 (discrete Fourier transform).
\leads The theorem is attributed to Napoleon Bonaparte, who was a good amateur geometer, though there is no evidence that he
found it; it was printed by W.~Rutherford in 1825. The inner Napoleon triangle reduces to a point exactly when
$\hat z_2=0$, that is when $ABC$ is itself equilateral and counterclockwise.
\end{solution}

\begin{solution}{pb:pdn}
\idea Erecting the same triangle on every side is a circulant map, diagonal on Fourier modes. The apex angle $\frac{2\pi k}n$
is tuned to kill exactly one mode. Since all these maps are diagonal in the same basis, they commute, and $n-2$ of them kill
$n-2$ modes, leaving the centroid (mode $0$) and a single regular polygon.

\writeup (a) Solving $z_k-p=e^{i\theta}(z_{k+1}-p)$: $p=\frac{z_k-e^{i\theta}z_{k+1}}{1-e^{i\theta}}$, so $A_\theta$ is linear. For
$z_k=\omega^{jk}$, $(A_\theta z)_k=\omega^{jk}\frac{1-e^{i\theta}\omega^j}{1-e^{i\theta}}$: the eigenvalue is
$\mu_j(\theta)=\frac{1-e^{i\theta}\omega^j}{1-e^{i\theta}}$. It is $1$ for $j=0$, and it is $0$ if and only if
$e^{i\theta}=\omega^{-j}$, that is $\theta=\frac{2\pi(n-j)}n$.

(b) $A_{2\pi k/n}$ kills the mode $j=n-k$. All the $A_\theta$ are diagonal in the Fourier basis, so they commute, and the
product of $A_{2\pi k/n}$ over $k=1,\dots,n-2$, in any order, multiplies the mode $j$ by $\prod_k\mu_j(\frac{2\pi k}n)$, which is
$0$ for $j=n-1,n-2,\dots,2$, $1$ for $j=0$, and some constant $c$ for $j=1$. The result is
$w_k=\hat z_0+c\,\hat z_1\omega^k$: a regular $n$-gon (a point if $c\hat z_1=0$) centred at the centroid $\hat z_0$ of $z$.

(c) $n=3$: one step, $\theta=120^\circ$, and this is Napoleon's theorem (Problem~\ref{pb:napoleon}). $n=4$: $\theta=90^\circ$ erects
squares on the sides and takes their centres; $\theta=180^\circ$ takes midpoints (the apex of a flat isosceles triangle is the
midpoint of its base). So: \emph{the centres of the squares erected on the sides of a quadrilateral form a quadrilateral whose
midpoint polygon is a square}, and by Varignon (Problem~\ref{pb:midpoint}) the order can be reversed: the centres of the
squares erected on the sides of the Varignon parallelogram form a square.
\tools Tome~4, Chapter~7 (discrete Fourier transform); Tome~0, Chapters~2 and~7 (rotations, diagonalisation).
\leads Petr (1908), Douglas (1940) and Neumann (1941) found this independently. For $n=4$ the statement contains van Aubel's
theorem: the segments joining the centres of opposite squares are equal and perpendicular.
\end{solution}

\begin{solution}{pb:rootwalk}
\idea Moments of $|S_m|$ expand into sums of products of steps, and $E[\xi^r\bar\xi^s]$ vanishes unless $n\mid r-s$: for small
moments and $n\ge3$, only the ``paired'' terms survive, as for a complex Gaussian. For a prime $p$, a return to $0$ is a
vanishing sum $\sum a_k\zeta^k$ with $a_k$ the number of steps in direction $\zeta^k$, which by Problem~\ref{pb:vanish} forces
$a_0=\dots=a_{p-1}$. The return probability is a multinomial coefficient, Stirling gives its size, and recurrence is decided
by the convergence of $\sum P(S_m=0)$.

\writeup (a) $\xi^r\bar\xi^s=\xi^{r-s}$, uniform over the powers $\zeta^{(r-s)k}$, so $E[\xi^r\bar\xi^s]=1$ if $n\mid r-s$ and $0$
otherwise. First, $E|S_m|^2=\sum_{a,b}E[\xi_a\bar\xi_b]=m$, since for $a\ne b$, $E[\xi_a\bar\xi_b]=E[\xi_a]E[\bar\xi_b]=0$.
Next, $E|S_m|^4=\sum_{a,b,c,d}E[\xi_a\xi_b\bar\xi_c\bar\xi_d]$. Group the factors by index: for each index, it appears $r$ times
unconjugated and $s$ times conjugated, with $r,s\le2$, and the term vanishes unless $n\mid r-s$ for every index. As
$|r-s|\le2<n$, this means $r=s$ for each index: the multisets $\{a,b\}$ and $\{c,d\}$ are equal. There are $m$ quadruples
with $a=b=c=d$ and $2m(m-1)$ with $a\ne b$ and $(c,d)\in\{(a,b),(b,a)\}$. So $E|S_m|^4=2m^2-m$.

(b) Let $a_k$ be the number of steps equal to $\zeta^k$ among the first $m$, so $S_m=\sum_ka_k\zeta^k$. By
Problem~\ref{pb:vanish}(a), $S_m=0$ if and only if $a_0=\dots=a_{p-1}$, which needs $p\mid m$. For $m=pN$ the vector
$(a_k)$ is multinomial, so $P(S_{pN}=0)=\frac{(pN)!}{N!^p}p^{-pN}$. Stirling's formula (Tome~4, Chapter~5) gives
$(pN)!\sim\sqrt{2\pi pN}\,(pN/e)^{pN}$ and $N!^p\sim(2\pi N)^{p/2}(N/e)^{pN}$, so the ratio is
$\sqrt{2\pi pN}\,(2\pi N)^{-p/2}=\sqrt p\,(2\pi N)^{-(p-1)/2}$.

(c) The walk lives in the group $\Z[\zeta]$, isomorphic to $\Z^{p-1}$ (basis $1,\zeta,\dots,\zeta^{p-2}$, by (a) of
Problem~\ref{pb:vanish}). The criterion of Tome~3, Chapter~7 (the walk returns infinitely often almost surely if
$\sum_mP(S_m=0)=\infty$, and finitely often if the series converges) has the same proof for any random walk on $\Z^d$. For
$p=3$, $P(S_{3N}=0)\sim\frac{\sqrt3}{2\pi N}$, a divergent series: recurrent. For $p\ge5$, the exponent $\frac{p-1}2\ge2$
makes the series converge: transient. Yet $S_m$ is a point of the plane. The reason is that $\Z[\zeta_5]$ is dense in $\C$:
the walker has room to wander in four dimensions, and hitting exactly $0$ is a four-dimensional event.
\tools Tome~3, Chapters~2 and~7 (moments, recurrence); Tome~4, Chapter~5 (Stirling); Problem~\ref{pb:vanish}.
\leads Part (a) is the start of the central limit theorem for $S_m/\sqrt m$: the limit is the complex Gaussian with
$E|Z|^2=1$, $E|Z|^4=2$, the same for all $n\ge3$. The return question, though, depends on arithmetic: for composite $n$ the
rank of $\Z[\zeta_n]$ is $\varphi(n)$, so the walk on the $n$-th roots of unity is recurrent exactly for $n\in\{3,4,6\}$, the
three lattices of the crystallographic restriction.
\end{solution}

\section*{Map of the chapter}

\begin{center}
\small
\adjustbox{max width=\linewidth}{\begin{tabular}{lll}
Problem & Topic & Tomes used\\\hline
\ref{pb:subsetsum} & subsets with sum divisible by $p$ & 0, 3\\
\ref{pb:dicemod} & dice sums modulo $m$ & 3, 0\\
\ref{pb:ramsum} & Ramanujan sums & 1, 0\\
\ref{pb:bricks} & tilings and bricks & 0, 4\\
\ref{pb:prouhet} & Prouhet--Thue--Morse & 3, 4, 0\\
\ref{pb:lightsout} & Lights Out on a ring & 2, 0\\
\ref{pb:sieving} & cyclic sieving & 2, 0, 4\\
\ref{pb:phipm1} & $\Phi_n(\pm1)$, cyclotomic units & 1, 0\\
\ref{pb:phipq} & coefficients of $\Phi_{pq}$ and $\Phi_{105}$ & 1, 3, 4\\
\ref{pb:vanish} & vanishing sums of roots of unity & 0, 2\\
\ref{pb:kronecker} & Kronecker's theorem & 0, 2\\
\ref{pb:primsum} & sums of primitive roots & 1, 2\\
\ref{pb:phimodp} & $\Phi_n$ modulo $p$ & 2, 1\\
\ref{pb:rattan} & rational tangents, lattice polygons & 2, 1\\
\ref{pb:ramacube} & Ramanujan's cube roots & 0, 2\\
\ref{pb:discphi} & discriminant of $\Phi_p$, Zolotarev & 0, 2, 1\\
\ref{pb:eisensine} & Eisenstein's sine proof & 0, 1\\
\ref{pb:jensen} & Poisson and Jensen & 0, 4\\
\ref{pb:midpoint} & midpoint polygons & 4, 0\\
\ref{pb:napoleon} & Napoleon's theorem & 0, 4\\
\ref{pb:pdn} & Petr--Douglas--Neumann & 4, 0\\
\ref{pb:rootwalk} & walks on roots of unity & 3, 4, 2\\
\end{tabular}}
\end{center}

\section*{Sources and further reading}

\begin{itemize}
\item H.~Weber, Leopold Kronecker, \emph{Jahresbericht der Deutschen Mathematiker-Vereinigung} 2 (1893) (the epigraph,
which Weber reports as a remark made by Kronecker in a lecture at the 1886 meeting of natural scientists in Berlin).
\item A.~Migotti, Zur Theorie der Kreisteilungsgleichung, \emph{Sitzungsberichte der Akademie der Wissenschaften Wien} 87
(1883), 7--14; J.~Suzuki, On coefficients of cyclotomic polynomials, \emph{Proceedings of the Japan Academy} 63 (1987),
279--280; E.~Lehmer, On the magnitude of the coefficients of the cyclotomic polynomial, \emph{Bulletin of the American
Mathematical Society} 42 (1936), 389--392 (which publishes Schur's argument from his letter to Landau); T.~Y.~Lam and K.~H.~Leung, On the cyclotomic polynomial $\Phi_{pq}(X)$, \emph{American Mathematical Monthly} 103
(1996), 562--564.
\item S.~Ramanujan, On certain trigonometrical sums and their applications in the theory of numbers, \emph{Transactions of the
Cambridge Philosophical Society} 22 (1918), 259--276.
\item N.~G.~de Bruijn, Filling boxes with bricks, \emph{American Mathematical Monthly} 76 (1969), 37--40; S.~Wagon, Fourteen
proofs of a result about tiling a rectangle, \emph{American Mathematical Monthly} 94 (1987), 601--617.
\item J.-P.~Allouche and J.~Shallit, The ubiquitous Prouhet--Thue--Morse sequence, in \emph{Sequences and Their Applications},
Springer, 1999, 1--16.
\item K.~Sutner, Linear cellular automata and the Garden-of-Eden, \emph{Mathematical Intelligencer} 11 (1989), 49--53.
\item V.~Reiner, D.~Stanton and D.~White, The cyclic sieving phenomenon, \emph{Journal of Combinatorial Theory A} 108 (2004),
17--50.
\item N.~G.~de Bruijn, On the factorization of cyclic groups, \emph{Indagationes Mathematicae} 15 (1953), 370--377; T.~Y.~Lam and
K.~H.~Leung, On vanishing sums of roots of unity, \emph{Journal of Algebra} 224 (2000), 91--109.
\item L.~Kronecker, Zwei Sätze über Gleichungen mit ganzzahligen Coefficienten, \emph{Journal für die reine und angewandte
Mathematik} 53 (1857), 173--175; D.~H.~Lehmer, Factorization of certain cyclotomic functions, \emph{Annals of Mathematics} 34
(1933), 461--479.
\item I.~Niven, \emph{Irrational Numbers}, Carus Mathematical Monographs 11, MAA, 1956.
\item B.~C.~Berndt, \emph{Ramanujan's Notebooks, Part IV}, Springer, 1994.
\item G.~Zolotarev, Nouvelle démonstration de la loi de réciprocité de Legendre, \emph{Nouvelles Annales de Mathématiques} (2) 11
(1872), 354--362; F.~Lemmermeyer, \emph{Reciprocity Laws: From Euler to Eisenstein}, Springer, 2000.
\item G.~Eisenstein, Applications de l'algèbre à l'arith\-métique trans\-cendante, \emph{Journal für die reine und angewandte
Mathematik} 29 (1845), 177--184.
\item J.~L.~W.~V.~Jensen, Sur un nouvel et important théorème de la théorie des fonctions, \emph{Acta Mathematica} 22 (1899),
359--364.
\item E.~R.~Berlekamp, E.~N.~Gilbert and F.~W.~Sinden, A polygon problem, \emph{American Mathematical Monthly} 72 (1965),
233--241.
\item H.~S.~M.~Coxeter and S.~L.~Greitzer, \emph{Geometry Revisited}, MAA, 1967 (Napoleon's theorem).
\item J.~Douglas, Geometry of polygons in the complex plane, \emph{Journal of Mathematics and Physics} 19 (1940), 93--130;
B.~H.~Neumann, Some remarks on polygons, \emph{Journal of the London Mathematical Society} 16 (1941), 230--245.
\end{itemize}
